\section{Statistical modelling ofCPU forward and backward passes}
\label{sec:chaining-runl-repeat-experiment}

The first loop in Algorithm~\ref{alg:chaining-simd-backtrack} can process
\(L\) anchors at once only when \(\operatorname{Run}_L(j,c)\) is true.  The
score and equal-score tie rule select the first chain end.  The predicate then
requires each of the next \(L\) predecessor links to move to the preceding
slot in the coordinate-sorted anchor array.

The \emph{first predecessor path} starts at this chain end and repeatedly
follows \(p_j\) until the next predecessor is \(-1\).  A \emph{retained
anchor} is a minimizer hit that remains in the coordinate-sorted anchor array
after minimap2 applies its occurrence-frequency filters.  A
\emph{repeat-derived anchor} is a retained anchor that links a selected
minimizer in one block to an occurrence in another block.  A
\emph{same-target anchor group} contains retained anchors with the same
reference-sequence identifier, target end coordinate, and orientation.  Its
members have different query end coordinates in these constructions.  The
first predecessor path contains at most one member of such a group.  The
remaining members still occupy array slots.  They can make \(p_j<j-1\), even
when consecutive path anchors are close in both sequences.

\subsubsection{Reference and query structures}

Figure~\ref{fig:chaining-runl-inputs} shows the three constructions used to
measure this effect.  The drawings show representative minimizer links and
are not to scale.

\begin{figure}[H]
\centering
\begin{tikzpicture}[
  x=1cm,y=1cm,font=\scriptsize,>=Latex,
  runl unit/.style={
    draw=black!60,minimum height=.48cm,align=center,
    inner sep=1.5pt,outer sep=0pt,rounded corners=0pt},
  runl unique/.style={
    runl unit,draw=chainGray!85!black,fill=chainGray!10},
  runl source/.style={
    runl unit,draw=chainBlue!90!black,fill=chainBlue!13},
  runl variant/.style={
    runl unit,draw=chainOrange!95!black,fill=chainOrange!16},
  runl selected/.style={
    runl unit,draw=chainGreen!80!black,fill=chainGreen!14,very thick},
  runl unused/.style={
    runl unit,draw=chainOrange!95!black,fill=chainOrange!10,dashed},
  runl cell/.style={
    runl unit,minimum width=.52cm,minimum height=.42cm,
    fill=chainGray!8,font=\tiny},
  runl chosen cell/.style={
    runl cell,draw=chainGreen!80!black,fill=chainGreen!16,very thick},
  runl main/.style={line width=.85pt,draw=chainGreen!80!black},
  runl extra/.style={
    line width=.75pt,draw=chainOrange!95!black,dashed},
  runl pred/.style={-{Latex[length=1.4mm]},thick,draw=black!72},
  runl note/.style={
    draw=black!45,fill=black!2,align=left,
    inner sep=3pt,rounded corners=0pt}
]
  \node[anchor=west,font=\bfseries] at (.10,9.25)
    {(a) Two internal blocks in an exact full-length mapping};

  \node[anchor=east] at (.72,8.55) {reference};
  \node[runl unique,minimum width=.70cm] (rUzero) at (1.25,8.55)
    {\(U_0\)};
  \node[runl source,minimum width=3.00cm] (rA) at (3.15,8.55)
    {\(A\)};
  \node[runl unique,minimum width=.70cm] (rUone) at (5.05,8.55)
    {\(U_1\)};
  \node[runl variant,minimum width=3.00cm] (rAprime) at (6.95,8.55)
    {\(A'_{\delta}\)};
  \node[runl unique,minimum width=.70cm] (rUtwo) at (8.85,8.55)
    {\(U_2\)};

  \node[anchor=east] at (.72,7.72) {query};
  \node[runl unique,minimum width=.70cm] (qUzero) at (1.25,7.72)
    {\(U_0\)};
  \node[runl source,minimum width=3.00cm] (qA) at (3.15,7.72)
    {\(A\)};
  \node[runl unique,minimum width=.70cm] (qUone) at (5.05,7.72)
    {\(U_1\)};
  \node[runl variant,minimum width=3.00cm] (qAprime) at (6.95,7.72)
    {\(A'_{\delta}\)};
  \node[runl unique,minimum width=.70cm] (qUtwo) at (8.85,7.72)
    {\(U_2\)};

  \draw[runl main]
    (rUzero.south)--(qUzero.north)
    (rA.south)--(qA.north)
    (rUone.south)--(qUone.north)
    (rAprime.south)--(qAprime.north)
    (rUtwo.south)--(qUtwo.north);
  \draw[runl extra] (rA.south)--(qAprime.north);
  \draw[runl extra] (rAprime.south)--(qA.north);

  \node[runl note,text width=3.55cm] at (11.60,8.46) {
    reference \(=\) query\\[2pt]
    \(\lvert A\rvert=32768\)\\
    \(\lvert U_i\rvert=512\)\\
    \(A\)-to-\(A'_{\delta}\) offset \(=33280\)
  };
  \draw[runl main] (9.72,7.60)--(10.38,7.60);
  \node[anchor=west] at (10.50,7.60) {selected equal-coordinate path};
  \draw[runl extra] (9.72,7.28)--(10.38,7.28);
  \node[anchor=west] at (10.50,7.28) {repeat-derived anchors};

  \draw[black!35] (.10,6.98)--(13.75,6.98);

  \node[anchor=west,font=\bfseries] at (.10,6.70)
    {(b) Exact spaced-copy control};

  \node[anchor=east] at (.72,6.05) {reference};
  \node[runl source,minimum width=1.50cm] at (1.48,6.05) {\(S\)};
  \node[runl unique,minimum width=1.50cm] at (3.00,6.05) {\(N^G\)};

  \node[anchor=east] at (.72,5.35) {query};
  \node[runl source,minimum width=.78cm] at (1.10,5.35) {\(S_0\)};
  \node[runl unique,minimum width=.72cm] at (1.87,5.35) {\(N^G\)};
  \node[runl source,minimum width=.78cm] at (2.64,5.35) {\(S_1\)};
  \node[runl unique,minimum width=.72cm] at (3.41,5.35) {\(N^G\)};
  \node at (4.13,5.35) {\(\cdots\)};
  \node[runl source,minimum width=1.02cm] at (5.05,5.35)
    {\(S_{R-1}\)};
  \node[runl unique,minimum width=.72cm] at (5.95,5.35) {\(N^G\)};

  \node[runl note,text width=1.55cm] at (7.30,5.70) {
    \(\lvert S\rvert=4096\)\\
    \(G=16384\)\\
    pitch \(=20480\)
  };

  \node[anchor=east] at (8.66,4.43) {\(r_t:\)};
  \node[runl cell] at (9.00,4.43) {\(0\)};
  \node[runl chosen cell] (previousV) at (9.58,4.43) {\(v\)};
  \node[runl cell] at (10.16,4.43) {\(R-1\)};

  \node[anchor=east] at (11.28,4.43) {\(r_{t+1}:\)};
  \node[runl cell] at (11.65,4.43) {\(0\)};
  \node[runl chosen cell] (currentV) at (12.23,4.43) {\(v\)};
  \node[runl cell] at (12.81,4.43) {\(R-1\)};

  \draw[runl pred]
    (currentV.north) .. controls +(0,.52) and +(0,.52) ..
    (previousV.north);
  \node[anchor=south,fill=white,inner sep=1pt] at (10.91,4.94)
    {\(p_j=j-R\neq j-1\quad(R\geq2)\)};
  \node[anchor=north] at (9.58,4.18) {index \(j-R\)};
  \node[anchor=north] at (12.23,4.18) {index \(j\)};

  \draw[black!35] (.10,3.55)--(13.75,3.55);

  \node[anchor=west,font=\bfseries] at (.10,3.28)
    {(c) Constructed coordinate control};

  \node[anchor=east] at (.72,2.70) {reference};
  \node[runl selected,minimum width=1.05cm] (rP) at (1.25,2.70)
    {\(P\ (40)\)};
  \node[runl unique,minimum width=1.05cm] (rG) at (2.45,2.70)
    {\(G_t\ (10)\)};
  \node[runl source,minimum width=.62cm] (rCzero) at (3.43,2.70)
    {\(c_0\)};
  \node[runl source,minimum width=.62cm] (rCone) at (4.18,2.70)
    {\(c_1\)};
  \node[runl variant,minimum width=.92cm] (rRone) at (5.06,2.70)
    {\(R_1\ (39)\)};
  \node[runl source,minimum width=.62cm] (rCtwo) at (5.94,2.70)
    {\(c_2\)};
  \node[runl variant,minimum width=.92cm] (rRtwo) at (6.82,2.70)
    {\(R_2\ (39)\)};
  \node at (7.62,2.70) {\(\cdots\)};
  \node[runl source,minimum width=.68cm] (rCm) at (8.42,2.70)
    {\(c_m\)};
  \node[runl variant,minimum width=.96cm] (rRm) at (9.36,2.70)
    {\(R_m\ (39)\)};

  \node[anchor=east] at (.72,2.03) {query};
  \node[runl selected,minimum width=1.05cm] (qP) at (1.25,2.03)
    {\(P\ (40)\)};
  \node[runl unique,minimum width=1.05cm] (qG) at (2.45,2.03)
    {\(G_q\ (10)\)};
  \node[runl source,minimum width=.62cm] (qCzero) at (3.43,2.03)
    {\(c_0\)};
  \node[runl source,minimum width=.62cm] (qCone) at (4.18,2.03)
    {\(c_1\)};
  \node[runl variant,minimum width=.92cm] (qRone) at (5.06,2.03)
    {\(R_1\ (39)\)};
  \node[runl source,minimum width=.62cm] (qCtwo) at (5.94,2.03)
    {\(c_2\)};
  \node[runl variant,minimum width=.92cm] (qRtwo) at (6.82,2.03)
    {\(R_2\ (39)\)};
  \node at (7.62,2.03) {\(\cdots\)};
  \node[runl source,minimum width=.68cm] (qCm) at (8.42,2.03)
    {\(c_m\)};
  \node[runl variant,minimum width=.96cm] (qRm) at (9.36,2.03)
    {\(R_m\ (39)\)};

  \draw[runl main]
    (rP.south)--(qP.north)
    (rG.south)--(qG.north)
    (rCzero.south)--(qCzero.north)
    (rCone.south)--(qCone.north)
    (rRone.south)--(qRone.north)
    (rCtwo.south)--(qCtwo.north)
    (rRtwo.south)--(qRtwo.north)
    (rCm.south)--(qCm.north)
    (rRm.south)--(qRm.north);

  \node[runl note,text width=3.35cm] at (11.84,2.36) {
    one isolated 25-base block per intended anchor\\
    \(N\)-filled gaps set positions
  };

  \node[anchor=east] at (.72,.82) {anchor array};
  \node[runl selected,minimum width=.82cm] (aP) at (1.18,.82) {\(P\)};
  \node[runl selected,minimum width=.82cm] (aG) at (2.08,.82) {\(G\)};
  \node[runl selected,minimum width=.58cm] (aCzero) at (2.82,.82)
    {\(c_0\)};
  \node[runl selected,minimum width=.58cm] (aCone) at (3.50,.82)
    {\(c_1\)};
  \node[runl unused,minimum width=.82cm] (aRone) at (4.30,.82)
    {\(R_1\)};
  \node[runl selected,minimum width=.58cm] (aCtwo) at (5.10,.82)
    {\(c_2\)};
  \node[runl unused,minimum width=.82cm] (aRtwo) at (5.90,.82)
    {\(R_2\)};
  \node at (6.62,.82) {\(\cdots\)};
  \node[runl selected,minimum width=.64cm] (aCm) at (7.35,.82)
    {\(c_m\)};
  \node[runl selected,minimum width=1.08cm] (aRm) at (8.22,.82)
    {\(e_m\in R_m\)};

  \draw[runl pred]
    (aRm.north) .. controls +(0,.28) and +(0,.28) .. (aCm.north);
  \draw[runl pred,dashed]
    (aCm.north) .. controls +(0,.60) and +(0,.60) .. (aCtwo.north);
  \draw[runl pred]
    (aCtwo.north) .. controls +(0,.60) and +(0,.60) .. (aCone.north);
  \draw[runl pred]
    (aCone.north) .. controls +(0,.28) and +(0,.28) .. (aCzero.north);
  \draw[runl pred]
    (aCzero.north) .. controls +(0,.28) and +(0,.28) .. (aG.north);
  \draw[runl pred]
    (aG.north) .. controls +(0,.28) and +(0,.28) .. (aP.north);

  \node[runl note,text width=4.35cm] at (11.28,.84) {
    first path:
    \(e_m\to\cdots\to r_{m,1}\to c_m\to\cdots\to c_0\to G\to P\)\\[2pt]
    \(R_1,\ldots,R_{m-1}\) stay in the array but are not on this path
  };
\end{tikzpicture}
\caption{Reference and query structures used to measure the first-path
\(\operatorname{Run}_L\) branch.  Green vertical lines connect matching
blocks at equal reference and query coordinates.  Green boxes in the
anchor-array row show the first predecessor path.  Dashed orange links in
panel~(a) are repeat-derived.  Dashed orange boxes in panel~(c)'s anchor-array
row are outside that path.  Panel~(a) keeps a full-length mapping while
placing \(A\) and \(A'\) inside the same sequence.  Panel~(b) isolates the
array-index effect of exact copies.  Panel~(c) constructs fixed anchor
coordinates from isolated 25-base blocks that each emit one minimizer.}
\label{fig:chaining-runl-inputs}
\end{figure}

The mosaic construction in panel~(a) uses a reference and query with identical
bytes:
\[
  T=Q=U_0\,A\,U_1\,A'_{\delta}\,U_2.
\]
Here \(\lvert A\rvert=32768\), every \(U_i\) has length 512, and the
complete sequence has length 67072.  The configured parameter \(\delta\)
sets the substitution fraction between \(A\) and \(A'_{\delta}\).  The
generator applies the exact count
\[
  s_{\delta}=
  \left\lfloor\delta\lvert A\rvert+\frac{1}{2}\right\rfloor.
\]
Every substituted base is changed to a different base.  The realized fraction
of positions at which the two blocks differ is therefore
\[
  \widehat{\delta}=\frac{s_{\delta}}{\lvert A\rvert}.
\]
This value is fixed when the input is generated.  The dense sweep uses
\[
  \delta_i=\frac{i}{400},
  \qquad i=0,\ldots,80.
\]
These 81 values run from 0 to 20 percent in steps of 0.25 percentage points.
Within each input generated from one fixed random-number-generator seed, the
substitutions for a smaller \(\delta\) form a subset of those used for every
larger value.  Both \(T\) and \(Q\) contain the same \(A\) and
\(A'_{\delta}\) blocks, so their query/reference error rate is zero.  The
quantities \(\delta\) and \(\widehat{\delta}\) describe only the internal
difference between the two blocks inside each sequence.

An independently generated \(A'\) provides a separate control.  It is
generated without copying or mutating \(A\).  Across the three inputs
generated from fixed random-number-generator seeds, no anchor on the first
predecessor path shared its reference-sequence identifier, target end
coordinate, and orientation with another retained anchor.  The path that
links equal target and query coordinates remains valid and has the maximum
score in every case.
Table~\ref{tab:chaining-runl-mosaic-divergence} lists the realized divergence
and corresponding first-path measurements.

The exact-copy control in panel~(b) uses
\[
  T=S\,N^G,
  \qquad
  Q=(S\,N^G)^R,
  \qquad
  \lvert S\rvert=4096,
  \qquad
  G=16384.
\]
Here \(R\) is the number of query copies.  The \(N\)-filled spacer blocks emit
no minimizers and make the copy pitch 20480.  In the generated inputs, each
target minimizer that passed occurrence filtering produced one anchor for
every query copy.  Let
\(j_{t,v}\) be the array index for copy \(v\), where \(0\leq v<R\).
Consecutive target minimizers satisfy
\[
  p_{j_{t+1,v}}=j_{t,v}=j_{t+1,v}-R.
\]
For \(R\geq2\), the first predecessor path contains no run of at least two unit
predecessor relations.  The construction therefore predicts zero vector
batches for every tested \(L\geq2\).

Panel~(c) is a constructed coordinate control.  Every intended anchor comes
from a distinct 25-base non-\(N\) block surrounded by \(N\) bases.  Each such
block emits one minimizer for \(k=15\) and \(w=10\).  This isolated block is
called a minimizer island.  The remaining positions contain \(N\).  This
control does not estimate the frequency of the structure in biological
sequence.  Write \(r_{k,\ell}\) for the \(\ell\)-th anchor in \(R_k\).  In the
equations below, \(c_k\) and \(r_{k,\ell}\) denote their target-query
coordinate pairs.  Let \(s_k=1\) for odd \(k\) and \(s_k=-1\) for even \(k\).
The following pairs give the target and query coordinate increments:
\[
\begin{aligned}
  c_k-c_{k-1}
    &=(3500,\ 3500-100s_k),\\
  r_{k,\ell}-c_k
    &=(500+38(\ell-1),\
       500-450s_k+38(\ell-1)),
       &&1\leq\ell\leq39.
\end{aligned}
\]
The \texttt{map-ont} preset sets the short-range chaining bandwidth to 500.
Every anchor in \(R_k\) and \(c_{k+1}\) differ by 550 between the target and
query distances, so those anchors fail the bandwidth check.  The preceding
\(c_k\) remains valid, and the recorded predecessor array confirms
\(p(c_{k+1})=c_k\).  The 39 anchors of \(R_k\) lie between \(c_k\) and
\(c_{k+1}\) in array order, which gives an index jump of 40.

Let \(h\) be the last anchor in \(P\), and let \(e_k=r_{k,39}\) be the last
anchor in \(R_k\).  With the unmodified \texttt{map-ont} scoring coefficients,
the construction gives
\[
  f(h)=600,
  \qquad
  f(c_k)=88,
  \qquad
  f(e_k)=615.
\]
All \(e_k\) values share the maximum score.  The experiment orders equal-score
ends by decreasing anchor-array index, matching BioseqDB's tie rule.  This
selects \(e_m\).  The first measured path passes through \(R_m\) and then
follows \(c_m,\ldots,c_0,G,P\).  The worst-case argument for all minimap2 cut
searches concerns the total work over every selected chain end and remains a
separate result.

\subsubsection{Metrics and results}

Minimap2 2.31-r1302 was run with the \texttt{map-ont} preset.  Filtering and
scoring parameters were not changed.  Instrumentation recorded the score and
predecessor arrays.  The only candidate-order change placed equal-score ends
in decreasing anchor-array index, matching BioseqDB's tie rule.  The tested
widths were \(L\in\{2,4,8,16,32,64\}\).  The branch counters were recomputed
from the recorded arrays.  For the mosaic and spaced-copy inputs, an
independent checker reconstructed the expected path coordinates and compared
them with the recorded predecessor path.  For the coordinate control, the
checker compared the recorded predecessor array with the array specified by
the construction.  The mosaic and spaced-copy inputs used three fixed
random-number-generator seeds.

Let \(q\) be the number of anchors on the first predecessor path.  Let
\(B_L\) be the number of evaluations for which
\(\operatorname{Run}_L\) is true, \(E_L\) the number of times the first-path
loop evaluates \(\operatorname{Run}_L\), and \(P_L\) the number that satisfy
both bounds \(j\geq L-1\) and
\(c\leq n-L\).  Three different ratios are useful:
\[
  \operatorname{Fire}_L=\frac{B_L}{E_L},
  \qquad
  \operatorname{Contiguous}_L=\frac{B_L}{P_L},
  \qquad
  \operatorname{Coverage}_L=\frac{LB_L}{q}.
\]
The first ratio is the observed branch-firing rate.  The second isolates the
predecessor test after the bounds pass.  The third is the fraction of path
anchors consumed by vector batches.  A successful evaluation consumes
\(L\) anchors, so firing rate and anchor coverage are not interchangeable.
An empty denominator gives an undefined value.

For maximal runs of unit predecessor relations with lengths \(r_i\), the
greedy loop gives
\[
  B_L=\sum_i\left\lfloor\frac{r_i}{L}\right\rfloor,
  \qquad
  E_L=q-(L-1)B_L.
\]
In the coordinate control, the path ends at anchor-array index zero, so
\(p_0=-1=0-1\) contributes one unit relation.  The \(m=1\) path is one run of
length 91.  For
\(m\geq2\), the only maximal runs have lengths 39 and 52.  Thus
\[
  B_L=
  \left\lfloor\frac{39}{L}\right\rfloor+
  \left\lfloor\frac{52}{L}\right\rfloor,
  \qquad
  q=m+90.
\]
For fixed \(L\) with \(B_L>0\), both \(\operatorname{Fire}_L\) and
\(\operatorname{Coverage}_L\) decrease as \(\Theta(1/m)\).  For \(L=64\),
both values are zero for every \(m\geq2\).

The row prefixes in the next table refer to
Figure~\ref{fig:chaining-runl-inputs}.  In the mosaic rows, \emph{exact}
means \(A'=A\), while \emph{independent} means that \(A'\) was generated
without copying \(A\).  The itemized notes below the table describe every
input family.

\begin{table}[H]
\centering
\caption{Observed first-path branch-firing rate \(100B_L/E_L\), in percent.
The mosaic rows and the \(R=1\) spaced-copy row report medians over three
inputs generated from fixed random-number-generator seeds.  The \(R\geq2\)
spaced-copy row covers the tested powers of two from 2 through 64, all of
which gave zero.  Coordinate-control values are deterministic and agreed with
the recorded counters.}
\label{tab:chaining-runl-firing}
\scriptsize
\setlength{\tabcolsep}{3.2pt}
\begin{tabular}{lrrrrrr}
\toprule
Input & \(L=2\) & \(L=4\) & \(L=8\) & \(L=16\) & \(L=32\) & \(L=64\) \\
\midrule
(a) Mosaic, independent \(A'\) & 100.00 & 99.94 & 99.74 & 99.24 & 97.01 & 91.12 \\
(a) Mosaic, exact \(A'=A\)     &   1.17 &  0.59 &  0.29 &  0.14 &  0.07 &  0.02 \\
(b) Spaced copies, \(R=1\)     & 100.00 & 98.97 & 94.12 & 88.89 & 80.00 & 66.67 \\
(b) Spaced copies, \(R\geq2\) &   0.00 &  0.00 &  0.00 &  0.00 &  0.00 &  0.00 \\
(c) Coordinate control, \(m=2\)    & 95.74 & 84.62 & 45.45 & 29.41 & 6.67 & 0.00 \\
(c) Coordinate control, \(m=1024\) &  4.21 &  2.10 &  0.96 &  0.48 & 0.19 & 0.00 \\
\bottomrule
\end{tabular}
\end{table}

The row prefixes refer to the panels in
Figure~\ref{fig:chaining-runl-inputs}.
\begin{itemize}
  \item \textbf{(a) Mosaic, exact.}
    The target and query are both
    \(U_0 A U_1 A' U_2\), and \(A'=A\).  The word \emph{mosaic} refers to
    this sequence layout.  The word \emph{exact} means that it contains two
    exact copies of \(A\).

  \item \textbf{(a) Mosaic, independent.}
    The layout is the same, but \(A'\) is generated without copying or
    mutating \(A\).  The same complete sequence is used as the target and
    query.  Its equal-coordinate path is therefore exact.

  \item \textbf{(b) Spaced copies.}
    The reference contains one source block \(S\).  The query contains \(R\)
    exact copies separated by \(N\)-filled gaps.  The \(R=1\) case is the
    baseline.

  \item \textbf{(c) Coordinate control.}
    Each minimizer island is an isolated 25-base non-\(N\) block that emits
    one minimizer.  These blocks and the surrounding \(N\)-filled gaps create
    the fixed anchor coordinates specified above.  The parameter \(m\) is the
    number of \(R_k\) groups.  This input checks the counters and does not
    represent a biological sequence distribution.
\end{itemize}

Table~\ref{tab:chaining-runl-mosaic-divergence} shows selected checkpoints
from the dense mosaic sweep.  Its notation and the percentage of path anchors
with same-target alternatives are defined below.
\begin{itemize}
  \item \textbf{Configured internal divergence \(\delta\).}
    This is the requested fraction of positions changed when constructing
    \(A'_{\delta}\) from \(A\).  It compares two blocks inside one sequence.
    Every mosaic input still has \(T=Q\).

  \item \textbf{Substitution count \(s_{\delta}\).}
    The generator changes exactly
    \[
      s_{\delta}
      =\left\lfloor
        \delta\lvert A\rvert+\frac{1}{2}
       \right\rfloor
    \]
    positions.  Every selected base is replaced by a different base.

  \item \textbf{Realized internal divergence \(\widehat{\delta}\).}
    For any generated \(A'\),
    \[
      \widehat{\delta}
      =\frac{\lvert\{i:A_i\neq A'_i\}\rvert}{\lvert A\rvert}.
    \]
    In the fixed-count sweep this equals
    \(s_{\delta}/\lvert A\rvert\).  The table reports
    \(100\widehat{\delta}\) as a percentage.

  \item \textbf{Independent \(A'\).}
    This block is generated without copying or mutating \(A\), so it has no
    configured \(\delta\) or \(s_{\delta}\).  Its realized divergence is
    measured directly with the formula above.

  \item \textbf{Path anchors with same-target alternatives.}
    The raw experiment data call this measurement \emph{grouped path}.  Two
    retained anchors belong to the same group when they have the same
    reference-sequence identifier, target end coordinate, and orientation.
    Let \(g_{\mathrm{path}}\) be the number of anchors on the first predecessor
    path whose group contains at least one other retained anchor.  If the path
    contains \(q\) anchors, the reported percentage is
    \(100g_{\mathrm{path}}/q\).  The other anchor occupies a different
    anchor-array slot and can separate consecutive path anchors.  This
    percentage does not count groups, broken predecessor relations, or
    \(\operatorname{Run}_L\) evaluations.
\end{itemize}

% Generated by runl_dense_divergence_audit.py. Do not edit by hand.

% manifest sha256: ba60e4b6e82dd32547c854f6c3d90273050a519a777c327c2cce4d1829c8fbca

% runner aggregate sha256: ed840411aee062d90a0b6cab2f98b13cf1acd5cd2e2e46911f8f46fd2ebde386

\def\RunLDenseGroupedMedianCoordinates{%
  (0,97.68786127) (0.25024414,93.50982752) (0.50048828,89.55342903) (0.75073242,85.89075158)%
  (1.00097656,82.23867526) (1.2512207,78.86191931) (1.50146484,75.55307974) (1.74865723,72.23425592)%
  (1.99890137,69.33440257) (2.24914551,66.26081384) (2.49938965,63.27348509) (2.74963379,60.77755511)%
  (2.99987793,58.44486934) (3.25012207,55.95664993) (3.50036621,53.5011551) (3.75061035,51.25976012)%
  (4.00085449,49.20379604) (4.25109863,47.00779179) (4.50134277,45.13295536) (4.74853516,43.1629393)%
  (4.9987793,41.4114621) (5.24902344,39.71540491) (5.49926758,38.11897106) (5.74951172,36.41321784)%
  (5.99975586,34.78051134) (6.25,33.48081409) (6.50024414,32.36381167) (6.75048828,31.08824948)%
  (7.00073242,29.6522859) (7.25097656,28.57602833) (7.5012207,27.22664736) (7.75146484,26.01769912)%
  (7.99865723,25.08043758) (8.24890137,24.08301158) (8.49914551,23.03547012) (8.74938965,21.84508968)%
  (8.99963379,21.01029601) (9.24987793,19.99517801) (9.50012207,19.16438135) (9.75036621,18.34124367)%
  (10.00061035,17.23111576) (10.25085449,16.6091077) (10.50109863,15.78143833) (10.75134277,15.31661845)%
  (10.99853516,14.6142136) (11.2487793,14.02409639) (11.49902344,13.47826087) (11.74926758,13.05781137)%
  (11.99951172,12.32326478) (12.24975586,11.81482077) (12.5,11.26760563) (12.75024414,10.72463768)%
  (13.00048828,10.19569944) (13.25073242,9.724343) (13.50097656,9.28175018) (13.7512207,8.81808673)%
  (14.00146484,8.51303508) (14.24865723,8.22622108) (14.49890137,7.71890327) (14.74914551,7.50221364)%
  (14.99938965,7.17676402) (15.24963379,6.80666184) (15.49987793,6.68971617) (15.75012207,6.40592306)%
  (16.00036621,6.20130131) (16.25061035,5.95028067) (16.50085449,5.64057367) (16.75109863,5.38547844)%
  (17.00134277,5.0669446) (17.24853516,4.79743281) (17.4987793,4.57023733) (17.74902344,4.36702256)%
  (17.99926758,4.22795595) (18.24951172,4.17637138) (18.49975586,4.0138075) (18.75,3.8060061)%
  (19.00024414,3.6009967) (19.25048828,3.45492528) (19.50073242,3.33921978) (19.75097656,3.16313423)%
  (20.0012207,3.03516942)%
}

\def\RunLDenseGroupedMinimumCoordinates{%
  (0,97.63691522) (0.25024414,93.40420391) (0.50048828,89.4715676) (0.75073242,85.67526312)%
  (1.00097656,82.18782581) (1.2512207,78.82730924) (1.50146484,75.14868992) (1.74865723,71.86522262)%
  (1.99890137,68.93344026) (2.24914551,66.12127045) (2.49938965,62.95702373) (2.74963379,60.29458854)%
  (2.99987793,57.76103854) (3.25012207,55.45585643) (3.50036621,53.03078897) (3.75061035,50.91375441)%
  (4.00085449,48.60442733) (4.25109863,46.78362573) (4.50134277,45.01241888) (4.74853516,43.13316764)%
  (4.9987793,41.33354653) (5.24902344,39.45632942) (5.49926758,37.8170761) (5.74951172,36.07868223)%
  (5.99975586,34.19664269) (6.25,32.78714731) (6.50024414,31.49568552) (6.75048828,30.37469042)%
  (7.00073242,29.09352059) (7.25097656,27.77644519) (7.5012207,26.54838967) (7.75146484,25.39987204)%
  (7.99865723,24.44053708) (8.24890137,23.2238472) (8.49914551,22.23644217) (8.74938965,20.96748358)%
  (8.99963379,20.18730489) (9.24987793,19.42152071) (9.50012207,18.64637565) (9.75036621,17.81019558)%
  (10.00061035,16.81373335) (10.25085449,16.33667335) (10.50109863,15.7359186) (10.75134277,15.19859065)%
  (10.99853516,14.49228792) (11.2487793,13.86709367) (11.49902344,13.30664532) (11.74926758,12.38796415)%
  (11.99951172,11.66026871) (12.24975586,11.19104716) (12.5,10.6314948) (12.75024414,10.19413598)%
  (13.00048828,9.86231188) (13.25073242,9.32169456) (13.50097656,8.90454837) (13.7512207,8.5563211)%
  (14.00146484,8.27531012) (14.24865723,7.84627702) (14.49890137,7.54626292) (14.74914551,7.04958744)%
  (14.99938965,6.85512933) (15.24963379,6.67573841) (15.49987793,6.4061499) (15.75012207,6.10283518)%
  (16.00036621,5.78108736) (16.25061035,5.44130591) (16.50085449,5.15076382) (16.75109863,4.87452898)%
  (17.00134277,4.66832438) (17.24853516,4.44194997) (17.4987793,4.32553669) (17.74902344,4.10289286)%
  (17.99926758,3.82522407) (18.24951172,3.70044053) (18.49975586,3.59913621) (18.75,3.37762126)%
  (19.00024414,3.26008789) (19.25048828,3.0499002) (19.50073242,2.92285577) (19.75097656,2.82161645)%
  (20.0012207,2.66093053)%
}

\def\RunLDenseGroupedMaximumCoordinates{%
  (0,97.68865265) (0.25024414,93.5249501) (0.50048828,89.67617826) (0.75073242,85.90101219)%
  (1.00097656,82.29057907) (1.2512207,78.93396977) (1.50146484,75.58549888) (1.74865723,72.65045835)%
  (1.99890137,69.42451777) (2.24914551,66.83144634) (2.49938965,64.08417025) (2.74963379,61.17572089)%
  (2.99987793,58.44770848) (3.25012207,56.19980726) (3.50036621,53.60907701) (3.75061035,51.3351933)%
  (4.00085449,49.44107314) (4.25109863,47.55289421) (4.50134277,45.3399695) (4.74853516,43.44334698)%
  (4.9987793,41.62859662) (5.24902344,39.99039462) (5.49926758,38.24496671) (5.74951172,36.65005626)%
  (5.99975586,35.3648757) (6.25,34.15924946) (6.50024414,32.8099968) (6.75048828,31.49883842)%
  (7.00073242,30.30691562) (7.25097656,28.81641157) (7.5012207,27.73945131) (7.75146484,26.81008188)%
  (7.99865723,25.99179207) (8.24890137,24.6879781) (8.49914551,23.56451613) (8.74938965,22.51410153)%
  (8.99963379,21.22409483) (9.24987793,20.11128135) (9.50012207,19.19549477) (9.75036621,18.3446273)%
  (10.00061035,17.63756656) (10.25085449,16.83870968) (10.50109863,15.97160603) (10.75134277,15.38958417)%
  (10.99853516,14.80406386) (11.2487793,14.2408377) (11.49902344,13.81589553) (11.74926758,13.38272003)%
  (11.99951172,12.89078197) (12.24975586,12.19473008) (12.5,11.83089535) (12.75024414,11.43361129)%
  (13.00048828,10.77700265) (13.25073242,10.29176018) (13.50097656,9.7169887) (13.7512207,9.46432861)%
  (14.00146484,9.00280336) (14.24865723,8.60714858) (14.49890137,8.28383368) (14.74914551,7.82741198)%
  (14.99938965,7.57442028) (15.24963379,7.27301919) (15.49987793,7.0195165) (15.75012207,6.6083888)%
  (16.00036621,6.27953544) (16.25061035,6.07663941) (16.50085449,5.6936648) (16.75109863,5.42188001)%
  (17.00134277,5.21125631) (17.24853516,4.82333146) (17.4987793,4.79013137) (17.74902344,4.61464509)%
  (17.99926758,4.47330447) (18.24951172,4.29659319) (18.49975586,4.15097364) (18.75,3.90806439)%
  (19.00024414,3.69688725) (19.25048828,3.59884837) (19.50073242,3.49302996) (19.75097656,3.47561464)%
  (20.0012207,3.36161357)%
}

\def\RunLDenseFireFourMedianCoordinates{%
  (0,0.58823529) (0.25024414,1.06277805) (0.50048828,1.5610652) (0.75073242,1.99795953)%
  (1.00097656,2.61177895) (1.2512207,3.19739276) (1.50146484,3.88271879) (1.74865723,4.4653349)%
  (1.99890137,5.14306744) (2.24914551,5.78218563) (2.49938965,6.65384615) (2.74963379,7.44263581)%
  (2.99987793,8.16346922) (3.25012207,8.87826706) (3.50036621,9.82496353) (3.75061035,10.61233153)%
  (4.00085449,11.46903038) (4.25109863,12.26238875) (4.50134277,13.05845743) (4.74853516,14.01020987)%
  (4.9987793,14.96756256) (5.24902344,16.17295748) (5.49926758,17.16571706) (5.74951172,18.29102167)%
  (5.99975586,19.33434959) (6.25,20.31330917) (6.50024414,21.13970588) (6.75048828,22.22371065)%
  (7.00073242,23.51407001) (7.25097656,24.46957007) (7.5012207,25.50774038) (7.75146484,26.7323862)%
  (7.99865723,27.86885246) (8.24890137,29.08570568) (8.49914551,30.25932177) (8.74938965,31.57400157)%
  (8.99963379,32.6959847) (9.24987793,33.95734983) (9.50012207,35.17904509) (9.75036621,36.44039176)%
  (10.00061035,37.98623064) (10.25085449,39.27571728) (10.50109863,40.49689441) (10.75134277,41.60464694)%
  (10.99853516,42.75722437) (11.2487793,43.88878522) (11.49902344,45.01705191) (11.74926758,45.94025278)%
  (11.99951172,46.876813) (12.24975586,48.01019408) (12.5,48.7433208) (12.75024414,49.69915764)%
  (13.00048828,51.15853659) (13.25073242,52.09531635) (13.50097656,53.36820084) (13.7512207,54.74064626)%
  (14.00146484,55.83886865) (14.24865723,57.0280252) (14.49890137,58.23594267) (14.74914551,59.87889661)%
  (14.99938965,60.48460145) (15.24963379,62.17150761) (15.49987793,62.98791822) (15.75012207,63.62359551)%
  (16.00036621,64.37485284) (16.25061035,65.24419156) (16.50085449,66.62662182) (16.75109863,67.49150898)%
  (17.00134277,69.3504569) (17.24853516,70.48975512) (17.4987793,71.06981416) (17.74902344,71.57655382)%
  (17.99926758,72.51148545) (18.24951172,73.05819021) (18.49975586,74.18436044) (18.75,75.25503531)%
  (19.00024414,75.77265584) (19.25048828,76.44879598) (19.50073242,76.99256111) (19.75097656,78.0101904)%
  (20.0012207,79.4133623)%
}

\def\RunLDenseFireFourMinimumCoordinates{%
  (0,0.58177647) (0.25024414,1.0083478) (0.50048828,1.47304988) (0.75073242,1.98446208)%
  (1.00097656,2.57267816) (1.2512207,3.19037657) (1.50146484,3.81693037) (1.74865723,4.37426264)%
  (1.99890137,5.0831793) (2.24914551,5.78061848) (2.49938965,6.58648339) (2.74963379,7.4092104)%
  (2.99987793,8.1622533) (3.25012207,8.80427118) (3.50036621,9.67112025) (3.75061035,10.40794979)%
  (4.00085449,11.24399359) (4.25109863,12.04216933) (4.50134277,13.01484541) (4.74853516,14.0038773)%
  (4.9987793,14.96622409) (5.24902344,16.02935954) (5.49926758,17.08206687) (5.74951172,17.97599901)%
  (5.99975586,19.11671924) (6.25,20.04365691) (6.50024414,21.03475307) (6.75048828,21.96969697)%
  (7.00073242,22.96968056) (7.25097656,24.32769615) (7.5012207,25.2007325) (7.75146484,25.90584879)%
  (7.99865723,26.82253848) (8.24890137,28.22304833) (8.49914551,29.42605527) (8.74938965,30.84083152)%
  (8.99963379,32.19720989) (9.24987793,33.85711971) (9.50012207,34.85845951) (9.75036621,36.26133692)%
  (10.00061035,37.63958083) (10.25085449,38.94037419) (10.50109863,40.44104571) (10.75134277,41.12367618)%
  (10.99853516,42.73143905) (11.2487793,43.39064522) (11.49902344,44.36936937) (11.74926758,45.153351)%
  (11.99951172,46.6461064) (12.24975586,47.49561489) (12.5,48.51778656) (12.75024414,49.53159259)%
  (13.00048828,50.84401058) (13.25073242,52.02716608) (13.50097656,53.30277141) (13.7512207,54.10628019)%
  (14.00146484,55.66058002) (14.24865723,56.92843159) (14.49890137,58.20269375) (14.74914551,59.36870383)%
  (14.99938965,60.40162455) (15.24963379,61.38229927) (15.49987793,61.79234472) (15.75012207,63.05658627)%
  (16.00036621,64.24384525) (16.25061035,65.02602934) (16.50085449,66.22754491) (16.75109863,67.18221363)%
  (17.00134277,68.29626008) (17.24853516,69.97765086) (17.4987793,70.15667744) (17.74902344,71.20603015)%
  (17.99926758,71.66666667) (18.24951172,72.44975833) (18.49975586,73.43429158) (18.75,74.86353002)%
  (19.00024414,75.42617362) (19.25048828,76.06299213) (19.50073242,76.41783171) (19.75097656,76.92715232)%
  (20.0012207,77.57656458)%
}

\def\RunLDenseFireFourMaximumCoordinates{%
  (0,0.59315837) (0.25024414,1.1842215) (0.50048828,1.78813559) (0.75073242,2.29263266)%
  (1.00097656,2.80365694) (1.2512207,3.31125828) (1.50146484,3.97625045) (1.74865723,4.71611722)%
  (1.99890137,5.29739777) (2.24914551,6.01154975) (2.49938965,6.78374988) (2.74963379,7.57441792)%
  (2.99987793,8.47236181) (3.25012207,9.27898402) (3.50036621,10.17161992) (3.75061035,11.07314469)%
  (4.00085449,11.95379754) (4.25109863,12.66998342) (4.50134277,13.54366197) (4.74853516,14.5605345)%
  (4.9987793,15.37830874) (5.24902344,16.44949193) (5.49926758,17.33934886) (5.74951172,18.37750404)%
  (5.99975586,19.35809957) (6.25,20.35272915) (6.50024414,21.65546329) (6.75048828,22.58107746)%
  (7.00073242,23.98463227) (7.25097656,25.45121263) (7.5012207,26.49168101) (7.75146484,27.57562473)%
  (7.99865723,28.76712329) (8.24890137,30.09428224) (8.49914551,31.38686131) (8.74938965,32.88782816)%
  (8.99963379,34.13804277) (9.24987793,35.10445797) (9.50012207,36.24811904) (9.75036621,37.40432046)%
  (10.00061035,38.72030518) (10.25085449,39.53040126) (10.50109863,40.55068836) (10.75134277,41.80505415)%
  (10.99853516,43.18558028) (11.2487793,44.44236877) (11.49902344,45.86265931) (11.74926758,47.79898714)%
  (11.99951172,49.3158834) (12.24975586,50.55320861) (12.5,52.09997951) (12.75024414,53.30149502)%
  (13.00048828,54.18242959) (13.25073242,55.35904539) (13.50097656,56.72868888) (13.7512207,57.68978342)%
  (14.00146484,58.36635843) (14.24865723,60.09874327) (14.49890137,61.02040816) (14.74914551,62.71929825)%
  (14.99938965,63.26293804) (15.24963379,63.91872957) (15.49987793,64.95867769) (15.75012207,65.83273767)%
  (16.00036621,67.29514141) (16.25061035,68.4171959) (16.50085449,69.52122409) (16.75109863,70.27161724)%
  (17.00134277,70.76653307) (17.24853516,71.94226386) (17.4987793,72.60692464) (17.74902344,73.84694666)%
  (17.99926758,75.70680628) (18.24951172,76.531151) (18.49975586,76.95157449) (18.75,77.6652452)%
  (19.00024414,78.3877879) (19.25048828,79.71838614) (19.50073242,80.33769063) (19.75097656,80.92896175)%
  (20.0012207,81.92837466)%
}

\def\RunLDenseCoverageFourMedianCoordinates{%
  (0,2.31213873) (0.25024414,4.11976048) (0.50048828,5.96491228) (0.75073242,7.53990535)%
  (1.00097656,9.68802631) (1.2512207,11.67014949) (1.50146484,13.91055228) (1.74865723,15.75129534)%
  (1.99890137,17.82241352) (2.24914551,19.70977516) (2.49938965,22.18659827) (2.74963379,24.33667335)%
  (2.99987793,26.23003451) (3.25012207,28.04368776) (3.50036621,30.35326306) (3.75061035,32.19834179)%
  (4.00085449,34.13221811) (4.25109863,35.85830187) (4.50134277,37.53094306) (4.74853516,39.45686901)%
  (4.9987793,41.31755676) (5.24902344,43.55798617) (5.49926758,45.32287749) (5.74951172,47.24132416)%
  (5.99975586,48.94677601) (6.25,50.48668651) (6.50024414,51.74353206) (6.75048828,53.33547661)%
  (7.00073242,55.15132003) (7.25097656,56.44369315) (7.5012207,57.80030574) (7.75146484,59.34030571)%
  (7.99865723,60.71428571) (8.24890137,62.12998713) (8.49914551,63.44406016) (8.74938965,64.85964771)%
  (8.99963379,66.02316602) (9.24987793,67.28489638) (9.50012207,68.46265527) (9.75036621,69.63536625)%
  (10.00061035,71.01600837) (10.25085449,72.12272107) (10.50109863,73.13516545) (10.75134277,74.02503028)%
  (10.99853516,74.92339945) (11.2487793,75.77929923) (11.49902344,76.60809286) (11.74926758,77.26870118)%
  (11.99951172,77.92333039) (12.24975586,78.69537275) (12.5,79.18341103) (12.75024414,79.80676329)%
  (13.00048828,80.73129661) (13.25073242,81.30811157) (13.50097656,82.07190541) (13.7512207,82.87070561)%
  (14.00146484,83.49219063) (14.24865723,84.14810066) (14.49890137,84.79691764) (14.74914551,85.652418)%
  (14.99938965,85.96025424) (15.24963379,86.797007) (15.49987793,87.19144488) (15.75012207,87.49396427)%
  (16.00036621,87.84641337) (16.25061035,88.24755491) (16.50085449,88.87108405) (16.75109863,89.25248637)%
  (17.00134277,90.0505091) (17.24853516,90.52547132) (17.4987793,90.76330981) (17.74902344,90.96893313)%
  (17.99926758,91.34313962) (18.24951172,91.55891093) (18.49975586,91.99646785) (18.75,92.40404689)%
  (19.00024414,92.59822357) (19.25048828,92.84910815) (19.50073242,93.04864344) (19.75097656,93.41682723)%
  (20.0012207,93.91360206)%
}

\def\RunLDenseCoverageFourMinimumCoordinates{%
  (0,2.28718692) (0.25024414,3.91496189) (0.50048828,5.64283424) (0.75073242,7.49183072)%
  (1.00097656,9.55337951) (1.2512207,11.64677804) (1.50146484,13.69906962) (1.74865723,15.46730846)%
  (1.99890137,17.64234162) (2.24914551,19.70522269) (2.49938965,21.99904352) (2.74963379,24.24725187)%
  (2.99987793,26.22689611) (3.25012207,27.85879353) (3.50036621,29.98486418) (3.75061035,31.72578717)%
  (4.00085449,33.63142766) (4.25109863,35.38522954) (4.50134277,37.44079634) (4.74853516,39.44431061)%
  (4.9987793,41.31500683) (5.24902344,43.29682629) (5.49926758,45.17684887) (5.74951172,46.71274715)%
  (5.99975586,48.59663192) (6.25,50.06815813) (6.50024414,51.58603012) (6.75048828,52.96803653)%
  (7.00073242,54.39538425) (7.25097656,56.25450757) (7.5012207,57.40413926) (7.75146484,58.30791459)%
  (7.99865723,59.45119498) (8.24890137,61.13213624) (8.49914551,62.51612903) (8.74938965,64.07735697)%
  (8.99963379,65.5108459) (9.24987793,67.18636985) (9.50012207,68.15768302) (9.75036621,69.47148258)%
  (10.00061035,70.71163466) (10.25085449,71.83870968) (10.50109863,73.08959421) (10.75134277,73.64191578)%
  (10.99853516,74.90359897) (11.2487793,75.40562249) (11.49902344,76.1352657) (11.74926758,76.70660127)%
  (11.99951172,77.76349614) (12.24975586,78.34753255) (12.5,79.03420523) (12.75024414,79.6985247)%
  (13.00048828,80.53475074) (13.25073242,81.26657558) (13.50097656,82.03319169) (13.7512207,82.50460405)%
  (14.00146484,83.39233988) (14.24865723,84.09383033) (14.49890137,84.77928761) (14.74914551,85.39000241)%
  (14.99938965,85.91831822) (15.24963379,86.40924781) (15.49987793,86.61151715) (15.75012207,87.2243163)%
  (16.00036621,87.78534241) (16.25061035,88.14755413) (16.50085449,88.69286287) (16.75109863,89.11684565)%
  (17.00134277,89.60153932) (17.24853516,90.31327618) (17.4987793,90.38769625) (17.74902344,90.81877904)%
  (17.99926758,91.00529101) (18.24951172,91.31863727) (18.49975586,91.70606619) (18.75,92.25594618)%
  (19.00024414,92.46845109) (19.25048828,92.70633397) (19.50073242,92.83768627) (19.75097656,93.02474574)%
  (20.0012207,93.26076517)%
}

\def\RunLDenseCoverageFourMaximumCoordinates{%
  (0,2.33115121) (0.25024414,4.57437384) (0.50048828,6.78838575) (0.75073242,8.58038081)%
  (1.00097656,10.34455064) (1.2512207,12.04819277) (1.50146484,14.20993409) (1.74865723,16.52627357)%
  (1.99890137,18.28388132) (2.24914551,20.37215271) (2.49938965,22.54650417) (2.74963379,24.68780019)%
  (2.99987793,27.02139595) (3.25012207,29.03380868) (3.50036621,31.17382938) (3.75061035,33.24783584)%
  (4.00085449,35.19409689) (4.25109863,36.72194184) (4.50134277,38.52255428) (4.74853516,40.5355568)%
  (4.9987793,42.09345195) (5.24902344,44.05667174) (5.49926758,45.62444854) (5.74951172,47.38530638)%
  (5.99975586,48.98481215) (6.25,50.54751818) (6.50024414,52.50878875) (6.75048828,53.84676839)%
  (7.00073242,55.79316949) (7.25097656,57.7276725) (7.5012207,59.0425957) (7.75146484,60.3646833)%
  (7.99865723,61.76470588) (8.24890137,63.2622073) (8.49914551,64.66165414) (8.74938965,66.21816434)%
  (8.99963379,67.4617786) (9.24987793,68.39195577) (9.50012207,69.45935122) (9.75036621,70.50336646)%
  (10.00061035,71.65089042) (10.25085449,72.33667335) (10.50109863,73.17899492) (10.75134277,74.18321589)%
  (10.99853516,75.25038058) (11.2487793,76.18895116) (11.49902344,77.21377102) (11.74926758,78.553137)%
  (11.99951172,79.55854127) (12.24975586,80.35171863) (12.5,81.31095124) (12.75024414,82.03243589)%
  (13.00048828,82.54883125) (13.25073242,83.22255145) (13.50097656,83.98462524) (13.7512207,84.50568819)%
  (14.00146484,84.86594638) (14.24865723,85.76461169) (14.49890137,86.22927181) (14.74914551,87.06240487)%
  (14.99938965,87.32281573) (15.24963379,87.63307452) (15.49987793,88.11659193) (15.75012207,88.51513695)%
  (16.00036621,89.16646649) (16.25061035,89.65351684) (16.50085449,90.12237063) (16.75109863,90.43534034)%
  (17.00134277,90.63928772) (17.24853516,91.11610006) (17.4987793,91.38096764) (17.74902344,91.86633544)%
  (17.99926758,92.57362356) (18.24951172,92.87945535) (18.49975586,93.03367192) (18.75,93.29278053)%
  (19.00024414,93.55173791) (19.25048828,94.01996008) (19.50073242,94.2341479) (19.75097656,94.43647378)%
  (20.0012207,94.77374124)%
}

\def\RunLDenseSelectedTableRows{%
  \(\delta=0\) & 0 & 0.000 & 97.69 [97.64, 97.69] & 0.59 [0.58, 0.59] & 2.31 [2.29, 2.33] \\%
  \(\delta=0.0025\) & 82 & 0.250 & 93.51 [93.40, 93.52] & 1.06 [1.01, 1.18] & 4.12 [3.91, 4.57] \\%
  \(\delta=0.0050\) & 164 & 0.500 & 89.55 [89.47, 89.68] & 1.56 [1.47, 1.79] & 5.96 [5.64, 6.79] \\%
  \(\delta=0.0100\) & 328 & 1.001 & 82.24 [82.19, 82.29] & 2.61 [2.57, 2.80] & 9.69 [9.55, 10.34] \\%
  \(\delta=0.0200\) & 655 & 1.999 & 69.33 [68.93, 69.42] & 5.14 [5.08, 5.30] & 17.82 [17.64, 18.28] \\%
  \(\delta=0.0500\) & 1638 & 4.999 & 41.41 [41.33, 41.63] & 14.97 [14.97, 15.38] & 41.32 [41.32, 42.09] \\%
  \(\delta=0.1000\) & 3277 & 10.001 & 17.23 [16.81, 17.64] & 37.99 [37.64, 38.72] & 71.02 [70.71, 71.65] \\%
  \(\delta=0.2000\) & 6554 & 20.001 & 3.04 [2.66, 3.36] & 79.41 [77.58, 81.93] & 93.91 [93.26, 94.77] \\%
}

\def\RunLDenseIndependentTableRow{%
  Independent \(A'\) & \text{n/a} & 75.043 & 0.00 [0.00, 0.00] & 99.94 [99.90, 100.00] & 99.98 [99.98, 100.00] \\%
}


\begin{table}[H]
\centering
\caption{Selected checkpoints from the 81-level dense mosaic sweep.  Each
entry in the last three columns gives the median followed by the observed
minimum and maximum over three inputs generated from fixed
random-number-generator seeds, in percent.  The firing column is
\(100\operatorname{Fire}_4=100B_4/E_4\).  The coverage column is
\(100\operatorname{Coverage}_4=100(4B_4)/q\).  The independent block is a
separate control.  Its realized divergence is measured directly from the two
generated blocks.  The third column gives its median, whose range was 74.930
to 75.186 percent.}
\label{tab:chaining-runl-mosaic-divergence}
\scriptsize
\setlength{\tabcolsep}{3.0pt}
\begin{tabular}{lrrlll}
\toprule
Case & \(s_{\delta}\) & \(100\widehat{\delta}\) &
\shortstack{Path anchors with\\same-target alternatives} &
\(100\operatorname{Fire}_4\) &
\(100\operatorname{Coverage}_4\) \\
& & (\%) & (\%) & (\%) & (\%) \\
\midrule
\RunLDenseSelectedTableRows
\midrule
\RunLDenseIndependentTableRow
\bottomrule
\end{tabular}
\end{table}

The repeated blocks affect \(\operatorname{Run}_L\) through anchor-array order.
When \(A\) and \(A'_{\delta}\) share a minimizer, one target minimizer can
produce anchors at both query copies.  The first predecessor path uses one of
them.  The other anchor remains in the coordinate-sorted array and can occupy
an array position between two consecutive path anchors.  This changes a unit
predecessor relation \(p_j=j-1\) into a larger index jump.  Increasing internal
divergence removes some of these shared minimizers and therefore removes some
of the extra array entries.

A local probability bound describes how substitutions can enable one
previously false evaluation.  Consider a loop state \((j,c)\) for which
\(\operatorname{Run}_L(j,c)\) is false when \(A'=A\).  Assume that this state
has \(h\) fixed, pairwise-disjoint blocking minimizer windows
\(W_1,\ldots,W_h\).  Each \(W_i\) is the base interval that supports one
repeat-derived minimizer occurrence whose retained anchor blocks this
specific \(\operatorname{Run}_L(j,c)\) evaluation.  The interval contains
\[
  b=k+w-1
\]
bases because it covers the \(w\) overlapping \(k\)-mers that select the
minimizer.  The evaluation can become true only if every \(W_i\) receives at
least one substitution.  If one window remains unchanged, its selected
minimizer occurrence and repeat-derived anchor remain available.

The generator ranks positions deterministically.  For this probability model,
let \(\pi\) be a uniform random permutation of the block positions.  Let
\(S_{\delta}\) be the first \(s_{\delta}\) positions under \(\pi\).  Thus
\(S_{\delta}\) is a uniformly selected subset of the
\(M=\lvert A\rvert\) block positions and
\(\widehat{\delta}=s_{\delta}/M\).  For one chosen base from every blocking
window, the probability that all \(h\) bases belong to \(S_{\delta}\) is
\[
  \frac{(s_{\delta})_h}{(M)_h},
  \qquad
  (x)_h=x(x-1)\cdots(x-h+1).
\]
There are \(b^h\) ways to choose one base from every window.  A union bound
therefore gives
\[
\begin{aligned}
  \Pr_{\pi}\!\left[
    \operatorname{Run}_L(j,c)\text{ changes from false to true}
  \right]
  &\leq
  \Pr_{\pi}\!\left[
    S_{\delta}\cap W_i\neq\varnothing
    \text{ for every }i
  \right] \\
  &\leq
  b^h\frac{(s_{\delta})_h}{(M)_h}
  \leq
  (b\widehat{\delta})^h .
\end{aligned}
\]
The probability is zero when \(s_{\delta}<h\).  For a lower bound, choose one
fixed base in each window.  The event that all \(h\) chosen bases are
substituted is contained in the event above and has probability
\[
  \frac{(s_{\delta})_h}{(M)_h}.
\]
For the asymptotic statement, consider a family of blocks with
\(M\to\infty\).  Keep \(b\) and \(h\) fixed, with the experimental values
\(k=15\), \(w=10\), and hence \(b=24\).  Require \(s_{\delta}\geq h\) and
\(\widehat{\delta}=s_{\delta}/M\to0\).  In this family, the probability that
every blocking window receives a substitution has order
\[
  \Theta\!\left(\widehat{\delta}^{\,h}\right).
\]
Since \(k\) is fixed here, this is equivalently
\[
  \Theta\!\left((k\widehat{\delta})^h\right).
\]

The last \(\Theta\) statement concerns the event that all blocking windows
receive substitutions.  It is not a two-sided bound on
\(\operatorname{Run}_L\).  Necessity gives only
\[
  \Pr_{\pi}\!\left[
    \operatorname{Run}_L(j,c)\text{ changes from false to true}
  \right]
  =
  O\!\left((k\widehat{\delta})^h\right).
\]
A matching lower bound would require substitutions in the \(h\) windows to
make the branch true.  Minimap2 does not provide that guarantee.  Minimizer
reselection, anchor order, chaining scores, predecessor choices, and the array
bounds can still prevent the run.  Outside the disjoint-window model, one
substitution can affect several overlapping windows.  The exponent can then
be smaller than the number of windows.

Figure~\ref{fig:chaining-runl-mosaic-divergence} reports the complete dense
sweep.  It uses all 81 measured divergence levels.

\begin{figure}[H]
\centering
\begin{tikzpicture}[
  font=\scriptsize,
  mosaic grouped/.style={
    draw=chainBlue!90!black,line width=.85pt},
  mosaic grouped range/.style={
    draw=chainBlue!45,line width=.35pt},
  mosaic fire/.style={
    draw=chainOrange!95!black,line width=.85pt,dashed},
  mosaic fire range/.style={
    draw=chainOrange!55,line width=.35pt,densely dashed}
]
\begin{scope}[x=.48cm,y=.045cm]
  \path[fill=black!3,draw=black!20]
    (21.2,0) rectangle (23.2,100);

  \foreach \yy in {0,25,50,75,100} {
    \draw[black!12] (0,\yy)--(20.5,\yy);
    \node[anchor=east] at (-.25,\yy) {\yy};
  }

  \draw[mosaic grouped range]
    plot coordinates {\RunLDenseGroupedMinimumCoordinates};
  \draw[mosaic grouped range]
    plot coordinates {\RunLDenseGroupedMaximumCoordinates};
  \draw[mosaic fire range]
    plot coordinates {\RunLDenseFireFourMinimumCoordinates};
  \draw[mosaic fire range]
    plot coordinates {\RunLDenseFireFourMaximumCoordinates};

  \draw[mosaic grouped]
    plot[mark=*,mark size=.42pt]
    coordinates {\RunLDenseGroupedMedianCoordinates};
  \draw[mosaic fire]
    plot[mark=x,mark size=.42pt]
    coordinates {\RunLDenseFireFourMedianCoordinates};

  \draw[->] (0,0)--(0,105);
  \draw[->] (0,0)--(20.8,0);
  \foreach \xx in {0,5,10,15,20} {
    \draw (\xx,-2)--(\xx,2);
    \node[anchor=north] at (\xx,-4.5) {\xx};
  }

  \node[rotate=90] at (-1.55,50) {reported value (\%)};
  \node at (10,-13)
    {realized internal divergence \(100\widehat{\delta}\) (\%)};

  \node[circle,draw=chainBlue!90!black,fill=white,
        minimum size=3.1pt,inner sep=0pt]
    at (22.2,0) {};
  \node[rectangle,draw=chainOrange!95!black,fill=white,
        minimum size=3.1pt,inner sep=0pt]
    at (22.2,99.94) {};
  \node[rotate=90,text=black!58] at (22.2,50)
    {independent \(A'\) control};
  \node[anchor=north,align=center] at (22.2,-4.5)
    {75.043\\outside axis};

  \draw[mosaic grouped] (.4,110)--(1.4,110);
  \node[anchor=west,align=left] at (1.65,110)
    {path anchors with\\same-target alternatives};
  \draw[mosaic fire] (12.0,110)--(13.0,110);
  \node[anchor=west] at (13.25,110) {\(\operatorname{Fire}_4\)};
\end{scope}
\end{tikzpicture}
\caption{Internal divergence, percentage of path anchors with same-target
alternatives, and \(L=4\) branch firing in the mosaic experiment.  All 81
configured levels from 0 to 20 percent are shown.  The small markers report
medians over three inputs generated from fixed random-number-generator seeds.
The lighter lines give the observed minimum and maximum.  They are descriptive
ranges, not confidence intervals.  The independent \(A'\) control is outside
the plotted divergence range and is shown in a separate column.}
\label{fig:chaining-runl-mosaic-divergence}
\end{figure}

The dense sweep contains 243 mappings.  It uses 81 nested divergence levels
for each of three inputs generated from fixed random-number-generator seeds.
Three additional mappings provide the independent \(A'\) controls.  Minimap2
emitted an alignment covering at least 99\% of both query and target for all
246 mappings.  Each mapping had exactly one anchor with the maximum dynamic
programming score.  This anchor fixed the start of the first predecessor path.
An independent checker reconstructed and matched every recorded predecessor
path.  All query minimizers survived query-side frequency filtering.  No
reference seed hits were removed by reference-occurrence filtering.  Each
retained minimizer value occurred once or twice in the reference, below the
cutoff of ten.

In the exact mosaic, the median percentage of path anchors with same-target
alternatives was 97.69 percent, and the median \(\operatorname{Fire}_4\) was
0.59 percent.  At 20 percent configured divergence, these values were 3.04
percent and 79.41 percent.  In the independent control, no anchor on the first
predecessor path had a same-target alternative.  Its median
\(\operatorname{Fire}_4\) was 99.94 percent.  At each of the 80 increases in
divergence, the alternative-anchor percentage decreased and
\(\operatorname{Fire}_4\) increased for all three inputs generated from fixed
random-number-generator seeds.  The medians followed the same directions.

The measurements support the array-order explanation.  More internal
divergence leaves fewer repeat-derived anchors interleaved with anchors on the
first predecessor path.  This restores more consecutive relations
\(p_j=j-1\), so \(\operatorname{Run}_4\) becomes true more often.  The result
is empirical and does not identify a universal power law.  Within each
input generated from one fixed random-number-generator seed, the 81 cases
share a cumulative mutation set.  They measure the effect of adding
substitutions to the same base sequence.  Statistical independence does not
apply to these 81 points.

For fixed \(L\), coverage is determined by firing rate:
\[
  \operatorname{Coverage}_L
  =\frac{L\operatorname{Fire}_L}
         {1+(L-1)\operatorname{Fire}_L}.
\]
This follows from \(E_L=q-(L-1)B_L\).  It explains why the coverage values in
Table~\ref{tab:chaining-runl-mosaic-divergence} rise faster than the firing
rate.

The spaced-copy control reached zero firing and zero anchor coverage at
\(R=2\), and both remained zero through \(R=64\).  In this construction,
the extra exact-copy anchors made \(\operatorname{Run}_L\) false at every
evaluated loop state.  These data do not establish the frequency of this
pattern in real reads.

Every tested coordinate-control case produced a path with defined counters.
The tested values were \(m=1\) and every power of two from 2 through 1024.
Every measured value agreed with the closed form above.  At \(m=1024\), anchor
coverage was 8.08 percent for \(L=2\) and 7.90 percent for \(L=4\).  It was
7.18 percent for both \(L=8\) and \(L=16\), 5.75 percent for \(L=32\), and
zero for \(L=64\).

\subsubsection{Random-sequence controls}

This experiment measures how often \(\operatorname{Run}_L\) is true for
anchors obtained from random DNA.  No anchor coordinates were predeclared.
Each target began with independently sampled A, C, G, and T bases with equal
probability.  A random source segment was then copied into \(N\)
non-overlapping target intervals.  The query was a randomly oriented copy of
that segment.  The mutation model used a substitution probability of 0.001.
The insertion and deletion tests each used probability 0.0005 at every source
base visited by the generator.  Candidate insertion and deletion lengths were sampled from a Gaussian
distribution with a mean and standard deviation of one base.  They were
rounded to integers and clipped at zero.  Zero-length draws caused no edit.
Deletions were clipped at the end of the oriented source.  Minimap2 used its default \(k=15\),
\(w=10\), occurrence filtering, and chaining scores.  Mapping used one thread.

The instrumentation reported the branch-firing percentage \(100B_L/E_L\) for
the first predecessor path from the initial dynamic-programming call.  Its
starting anchor was selected by score and the equal-score tie rule.  The
experiment tested \(L\in\{4,8,16\}\).  Equal scores were ordered by decreasing
anchor-array index to match BioseqDB.  A value is defined only when minimap2
found a score-qualified chain end and the measured path had \(E_L>0\).  Runs
without such a path are omitted because no firing percentage exists when
\(E_L=0\).  A displayed zero means that a measured path existed and
\(B_L=0\).  Neither experiment runner saved minimap2's PAF output.  These
results describe the internal predecessor path.  Whether minimap2 later
emitted each path as an alignment is unknown.

\begin{figure}[H]
\centering
\begin{tikzpicture}[
  font=\scriptsize,
  runl four/.style={draw=chainBlue!90!black,line width=.85pt},
  runl eight/.style={
    draw=chainOrange!95!black,line width=.85pt,dashed},
  runl sixteen/.style={
    draw=chainGreen!70!black,line width=.9pt,densely dotted},
  runl mark four/.style={
    circle,draw=chainBlue!90!black,fill=white,
    minimum size=4.4pt,inner sep=0pt},
  runl mark eight/.style={
    rectangle,draw=chainOrange!95!black,fill=white,
    minimum size=4.4pt,inner sep=0pt},
  runl mark sixteen/.style={
    circle,draw=chainGreen!70!black,fill=chainGreen!55,
    minimum size=4.4pt,inner sep=0pt}
]
  % Panel (a): original 10 Mb / 1 kb occurrence sweep.
  \begin{scope}[x=1cm,y=.038cm,xshift=1.55cm,yshift=6.05cm]
    \node[anchor=west,font=\bfseries\footnotesize] at (-1.25,112)
      {(a) Original: \(10\,\mathrm{Mb}\) target,
       \(1\,\mathrm{kb}\) query, one input};

    \foreach \yy in {0,25,50,75,100} {
      \draw[black!12] (0,\yy)--(11.15,\yy);
      \node[anchor=east] at (-.13,\yy) {\yy};
    }
    \draw[->] (0,0)--(6.08,0);
    \draw[->] (6.72,0)--(11.30,0)
      node[anchor=west] {\(N\)};
    \draw[->] (0,0)--(0,105);
    \node[rotate=90] at (-1.02,50)
      {branch firing (\%)};
    \draw (6.16,-3.6)--(6.31,3.6);
    \draw (6.45,-3.6)--(6.60,3.6);

    \foreach \xx/\lab in {
      .50/1,1.7041/2,2.4085/3,2.9082/4,
      3.2959/5,4.50/10,10.65/5000} {
      \draw (\xx,-2.2)--(\xx,2.2);
      \node[anchor=north] at (\xx,-5.0) {\lab};
    }

    \draw[runl four] plot coordinates {
      (.50,97.9167) (1.7041,97.9167) (2.4085,97.9167)
      (2.9082,97.9167) (3.2959,97.9167) (4.50,92.0000)};
    \draw[runl eight] plot coordinates {
      (.50,82.1429) (1.7041,82.1429) (2.4085,82.1429)
      (2.9082,82.1429) (3.2959,82.1429) (4.50,85.1852)};
    \draw[runl sixteen] plot coordinates {
      (.50,45.8333) (1.7041,45.8333) (2.4085,45.8333)
      (2.9082,45.8333) (3.2959,45.8333) (4.50,47.8261)};

    \foreach \xx/\yy in {
      .50/97.9167,1.7041/97.9167,2.4085/97.9167,
      2.9082/97.9167,3.2959/97.9167,4.50/92.0000}
      \node[runl mark four] at (\xx,\yy) {};
    \foreach \xx/\yy in {
      .50/82.1429,1.7041/82.1429,2.4085/82.1429,
      2.9082/82.1429,3.2959/82.1429,4.50/85.1852}
      \node[runl mark eight] at (\xx,\yy) {};
    \foreach \xx/\yy in {
      .50/45.8333,1.7041/45.8333,2.4085/45.8333,
      2.9082/45.8333,3.2959/45.8333,4.50/47.8261}
      \node[runl mark sixteen] at (\xx,\yy) {};

    \node[runl mark four] at (10.65,97.9167) {};
    \node[runl mark eight] at (10.65,82.1429) {};
    \node[runl mark sixteen] at (10.65,45.8333) {};

    \draw[runl four] (5.65,108)--(6.15,108);
    \node[anchor=west] at (6.23,108) {\(L=4\)};
    \draw[runl eight] (7.35,108)--(7.85,108);
    \node[anchor=west] at (7.93,108) {\(L=8\)};
    \draw[runl sixteen] (9.05,108)--(9.55,108);
    \node[anchor=west] at (9.63,108) {\(L=16\)};
  \end{scope}

  % Panel (b): corrected ten-input compact series.
  \begin{scope}[x=1cm,y=.038cm,xshift=1.55cm,yshift=.70cm]
    \node[anchor=west,font=\bfseries\footnotesize] at (-1.25,112)
      {(b) Ten-input compact series: \(\lvert T\rvert=4.5\,\mathrm{Mb}\),
       \(\lvert Q\rvert=997\) to \(1002\,\mathrm{bp}\)};

    \path[fill=chainBlue!45,fill opacity=.15]
      (.20,93.1008) rectangle (11.7051,95.7639);
    \path[fill=chainOrange!55,fill opacity=.13]
      (.20,83.5000) rectangle (11.7051,88.3462);
    \path[fill=chainGreen!45,fill opacity=.13]
      (.20,45.9957) rectangle (11.7051,66.7293);

    \foreach \yy in {0,25,50,75,100} {
      \draw[black!12] (0,\yy)--(11.90,\yy);
      \node[anchor=east] at (-.13,\yy) {\yy};
    }
    \draw[->] (0,0)--(12.05,0)
      node[anchor=west] {\(N\)};
    \draw[->] (0,0)--(0,105);
    \node[rotate=90] at (-1.02,50)
      {branch firing (\%)};

    \foreach \xx in {
      .20,1.0805,1.7051,2.1897,3.6949,5.2000,6.0805,
      6.7051,7.1897,8.6949,10.2000,11.0805,11.7051}
      \draw (\xx,-2.0)--(\xx,2.0);
    \foreach \xx/\lab in {
      .20/20,2.1897/50,3.6949/100,
      7.1897/500,8.6949/1000,11.7051/4000}
      \node[anchor=north] at (\xx,-5.0) {\lab};

    \draw[runl four] (.20,93.7473)--(11.7051,93.7473);
    \draw[runl eight] (.20,87.2283)--(11.7051,87.2283);
    \draw[runl sixteen] (.20,51.1905)--(11.7051,51.1905);

    \foreach \xx in {
      .20,1.0805,1.7051,2.1897,3.6949,5.2000,6.0805,
      6.7051,7.1897,8.6949,10.2000,11.0805,11.7051} {
      \node[runl mark four] at (\xx,93.7473) {};
      \node[runl mark eight] at (\xx,87.2283) {};
      \node[runl mark sixteen] at (\xx,51.1905) {};
    }
    \node[anchor=east,text=black!60] at (11.75,65)
      {shading: middle 50\% of inputs};
  \end{scope}
\end{tikzpicture}
\caption{First-path branch firing on random-sequence occurrence sweeps.
The plotted value is \(100B_L/E_L\).  Horizontal positions are logarithmic in
\(N\) within each unbroken axis segment.  Panel~(a) shows every defined result
from the original sweep over \(N\) with \(k=15\) and \(w=10\).
The axis break spans the tested values between \(N=20\) and \(N=4000\), for
which no score-qualified chain end was available.  The \(N=5000\) point is
plotted separately.  Panel~(b) shows the median over ten generated inputs at
each tested \(N\).  Its bands span the middle half of those ten values.  Every
panel~(b) run had a defined first path.}
\label{fig:chaining-runl-random-occurrences}
\end{figure}

Panel~(a) records the original run.  Its missing middle prevents an estimate
of the trend over \(N\).  The ten-input compact series produced a defined value for
all 13 selected \(N\) values spanning 20 to 4000.  Its median firing rates were
93.75 percent for \(L=4\), 87.23 percent for \(L=8\), and 51.19 percent for
\(L=16\).  For each generated input, all \(N\) values used the same sequence
background.  The measured path length and branch counters stayed constant as
\(N\) changed.  The selected copy and absolute anchor-array positions could
change.  The recorded \(\operatorname{Run}_L\) outcomes produced the flat
lines.  The earlier three-input compact pilot is not plotted.

The size controls kept \(k\), \(w\), mutation probabilities, and minimap2
parameters fixed.  They tested \(N=20\) and \(N=4000\).
Figure~\ref{fig:chaining-runl-random-sizes} shows the three target/query size
pairs that produced defined percentages: approximately
\(4.5\,\mathrm{Mb}/1\,\mathrm{kb}\), \(40\,\mathrm{Mb}/4\,\mathrm{kb}\),
and \(80\,\mathrm{Mb}/8\,\mathrm{kb}\).  The
\(10\,\mathrm{Mb}/1\,\mathrm{kb}\), \(20\,\mathrm{Mb}/2\,\mathrm{kb}\),
and \(40\,\mathrm{Mb}/1\,\mathrm{kb}\) controls produced no defined
percentage at either tested \(N\).

\begin{figure}[H]
\centering
\begin{tikzpicture}[
  font=\scriptsize,
  runl mark four/.style={
    circle,draw=chainBlue!90!black,fill=white,line width=.8pt,
    minimum size=5.2pt,inner sep=0pt},
  runl mark eight/.style={
    rectangle,draw=chainOrange!95!black,fill=white,line width=.8pt,
    minimum size=5.2pt,inner sep=0pt},
  runl mark sixteen/.style={
    circle,draw=chainGreen!70!black,fill=chainGreen!55,line width=.8pt,
    minimum size=5.2pt,inner sep=0pt},
  runl whisk four/.style={draw=chainBlue!90!black,line width=.55pt},
  runl whisk eight/.style={draw=chainOrange!95!black,line width=.55pt},
  runl whisk sixteen/.style={draw=chainGreen!70!black,line width=.55pt}
]
  \node[runl mark four] at (5.05,5.55) {};
  \node[anchor=west] at (5.18,5.55) {\(L=4\)};
  \node[runl mark eight] at (6.42,5.55) {};
  \node[anchor=west] at (6.55,5.55) {\(L=8\)};
  \node[runl mark sixteen] at (7.79,5.55) {};
  \node[anchor=west] at (7.92,5.55) {\(L=16\)};

  \begin{scope}[x=1cm,y=.041cm,xshift=1.35cm,yshift=.62cm]
    \node[font=\bfseries] at (2.50,109) {(a) \(N=20\)};
    \foreach \yy in {0,25,50,75,100} {
      \draw[black!12] (0,\yy)--(5.00,\yy);
      \node[anchor=east] at (-.12,\yy) {\yy};
    }
    \draw[->] (0,0)--(5.15,0);
    \draw[->] (0,0)--(0,104);
    \node[rotate=90] at (-.92,50) {branch firing (\%)};

    \foreach \xx in {.85,2.50,4.15}
      \draw[black!28] (\xx,0)--(\xx,100);
    \foreach \xx/\lo/\hi in {
      .67/90.9091/97.7273,3.97/42.8571/80.0000} {
      \draw[runl whisk four] (\xx,\lo)--(\xx,\hi);
      \draw[runl whisk four] (\xx-.06,\lo)--(\xx+.06,\lo);
      \draw[runl whisk four] (\xx-.06,\hi)--(\xx+.06,\hi);
    }
    \foreach \xx/\lo/\hi in {
      .85/77.7778/95.4545,4.15/11.1111/66.6667} {
      \draw[runl whisk eight] (\xx,\lo)--(\xx,\hi);
      \draw[runl whisk eight] (\xx-.06,\lo)--(\xx+.06,\lo);
      \draw[runl whisk eight] (\xx-.06,\hi)--(\xx+.06,\hi);
    }
    \foreach \xx/\lo/\hi in {
      1.03/41.6667/78.5714,4.33/0.0000/50.0000} {
      \draw[runl whisk sixteen] (\xx,\lo)--(\xx,\hi);
      \draw[runl whisk sixteen] (\xx-.06,\lo)--(\xx+.06,\lo);
      \draw[runl whisk sixteen] (\xx-.06,\hi)--(\xx+.06,\hi);
    }
    \node[runl mark four] at (.67,93.7473) {};
    \node[runl mark eight] at (.85,87.2283) {};
    \node[runl mark sixteen] at (1.03,51.1905) {};
    \node[runl mark four] at (2.32,20.0000) {};
    \node[runl mark eight] at (2.50,0.0000) {};
    \node[runl mark sixteen] at (2.68,0.0000) {};
    \node[runl mark four] at (3.97,50.0000) {};
    \node[runl mark eight] at (4.15,12.5000) {};
    \node[runl mark sixteen] at (4.33,0.0000) {};

    \node[anchor=north,align=center] at (.85,-6)
      {4.5/1\\[-1pt]\tiny 10 defined};
    \node[anchor=north,align=center] at (2.50,-6)
      {40/4\\[-1pt]\tiny 2 defined};
    \node[anchor=north,align=center] at (4.15,-6)
      {80/8\\[-1pt]\tiny 3 defined};
  \end{scope}

  \begin{scope}[x=1cm,y=.041cm,xshift=7.65cm,yshift=.62cm]
    \node[font=\bfseries] at (2.50,109) {(b) \(N=4000\)};
    \foreach \yy in {0,25,50,75,100}
      \draw[black!12] (0,\yy)--(5.00,\yy);
    \draw[->] (0,0)--(5.15,0);
    \draw[->] (0,0)--(0,104);

    \foreach \xx in {.85,2.50,4.15}
      \draw[black!28] (\xx,0)--(\xx,100);
    \foreach \xx/\lo/\hi in {
      .67/90.9091/97.7273,3.97/42.8571/80.0000} {
      \draw[runl whisk four] (\xx,\lo)--(\xx,\hi);
      \draw[runl whisk four] (\xx-.06,\lo)--(\xx+.06,\lo);
      \draw[runl whisk four] (\xx-.06,\hi)--(\xx+.06,\hi);
    }
    \foreach \xx/\lo/\hi in {
      .85/77.7778/95.4545,4.15/11.1111/66.6667} {
      \draw[runl whisk eight] (\xx,\lo)--(\xx,\hi);
      \draw[runl whisk eight] (\xx-.06,\lo)--(\xx+.06,\lo);
      \draw[runl whisk eight] (\xx-.06,\hi)--(\xx+.06,\hi);
    }
    \foreach \xx/\lo/\hi in {
      1.03/41.6667/78.5714,4.33/0.0000/50.0000} {
      \draw[runl whisk sixteen] (\xx,\lo)--(\xx,\hi);
      \draw[runl whisk sixteen] (\xx-.06,\lo)--(\xx+.06,\lo);
      \draw[runl whisk sixteen] (\xx-.06,\hi)--(\xx+.06,\hi);
    }
    \node[runl mark four] at (.67,93.7473) {};
    \node[runl mark eight] at (.85,87.2283) {};
    \node[runl mark sixteen] at (1.03,51.1905) {};
    \node[runl mark four] at (2.32,20.0000) {};
    \node[runl mark eight] at (2.50,0.0000) {};
    \node[runl mark sixteen] at (2.68,0.0000) {};
    \node[runl mark four] at (3.97,80.0000) {};
    \node[runl mark eight] at (4.15,66.6667) {};
    \node[runl mark sixteen] at (4.33,50.0000) {};

    \node[anchor=north,align=center] at (.85,-6)
      {4.5/1\\[-1pt]\tiny 10 defined};
    \node[anchor=north,align=center] at (2.50,-6)
      {40/4\\[-1pt]\tiny 3 defined};
    \node[anchor=north,align=center] at (4.15,-6)
      {80/8\\[-1pt]\tiny 3 defined};
  \end{scope}

\end{tikzpicture}
\caption{Median first-path branch firing for random-sequence size controls
with a defined path.  Each point is the median of the defined per-input
percentages.  The number of contributing inputs is printed below it.  The
horizontal labels give target length in megabases followed by query length in
kilobases.  Query lengths were 997 to 1002 bases, 3998 to 4004 bases, and 7999
to 8008 bases.  Vertical whiskers span the smallest and largest per-input
percentages.  Within each size pair, markers and whiskers are offset
horizontally only for visibility.  A plotted zero
means that a measured path existed and \(B_L=0\).  Runs with no score-qualified
chain end are omitted.  At \(40\,\mathrm{Mb}/4\,\mathrm{kb}\) and \(N=20\),
one of three generated inputs had no defined percentage, so the median uses
the other two.  The omitted target/query size pairs are
\(10\,\mathrm{Mb}/1\,\mathrm{kb}\), \(20\,\mathrm{Mb}/2\,\mathrm{kb}\),
and \(40\,\mathrm{Mb}/1\,\mathrm{kb}\) at both tested \(N\) values.}
\label{fig:chaining-runl-random-sizes}
\end{figure}

In the 40-Mb/4-kb and 80-Mb/8-kb controls, default occurrence filtering left
short first paths.  Query-side filtering removed none of their generated
minimizers.  For both sizes, the adaptive reference cutoff
\texttt{mid\_occ} remained 10.  Most repeated source minimizers had
approximately \(N\) reference occurrences.  A few also matched elsewhere in
the random target.  These counts exceeded \texttt{mid\_occ}.  Reference
occurrence filtering removed between 98.79 and 98.90 percent of the raw anchor
hits.  For each consecutive group above the cutoff, \texttt{occ\_dist}=500
set the number of minimizers that could remain from its query-coordinate
span.  It kept those with the smallest reference occurrence counts.  Every
defined 4-kb input had \(q=8\), while the 8-kb inputs had
\(q\in\{15,16,17\}\).  The longer query produced a larger limit under this
rule.  Target length also affects the adaptive cutoff, so this comparison does
not isolate query length.

The 4.5-Mb/1-kb series used larger adaptive cutoffs.  They were 21 for
\(N=20\) and 4001 for \(N=4000\).  Reference occurrence filtering retained
every anchor hit.  On each validated first path, every predecessor index was
one less than the current index.  These paths contained 163 to 187 anchors and
gave the higher firing percentages shown in
Figure~\ref{fig:chaining-runl-random-sizes}.

After filtering, \(\operatorname{Run}_L\) is evaluated from the predecessor
indices on these short paths.  All recorded bounds tests passed.  A displayed
zero therefore means that the predecessor-index equalities in
\(\operatorname{Run}_L\) failed at every loop evaluation.  The observed
first-path profiles were
\[
\begin{array}{c|ccc}
q & \operatorname{Fire}_4 & \operatorname{Fire}_8 &
    \operatorname{Fire}_{16} \\
\hline
8  & 1/5=20\% & 0/8=0\%     & 0/8=0\% \\
15 & 3/6=50\% & 1/8=12.5\%  & 0/15=0\% \\
16 & 3/7=42.86\% & 1/9=11.11\% & 0/16=0\% \\
17 & 4/5=80\% & 2/3=66.67\% & 1/2=50\%
\end{array}
\]
Every defined 40-Mb/4-kb input had the first profile.  It provided one
four-wide batch and no eight-wide or sixteen-wide batch.  For the 80-Mb/8-kb
controls, the path lengths at \(N=20\) were 16, 15, and 17.  At \(N=4000\)
they were 16, 17, and 17.  A 17-anchor path was the only observed profile with
one sixteen-wide batch.  That batch was followed by one scalar evaluation, so
\(B_{16}=1\), \(E_{16}=2\), and the firing percentage was 50 percent.  The
other profiles gave zero.  These discrete outcomes produce the whisker from
0 to 50 percent in Figure~\ref{fig:chaining-runl-random-sizes}.

The 40-Mb and 80-Mb conditions contained only two or three inputs with a
defined percentage.  Their whiskers are descriptive ranges.  The results do
not establish a length trend.  Minimap2 required a chaining score of at least
40.  The omitted 40-Mb/4-kb input at \(N=20\) had no anchor reaching that
threshold.

These measurements describe opportunities for the first-path vector branch.
They do not measure BioseqDB execution time.  The cost of failed checks and
the speed of successful vector batches must be evaluated by a BioseqDB
runtime benchmark.  Initial dynamic-programming and later RMQ calls were
recorded separately.  Table~\ref{tab:chaining-runl-firing} and
Figures~\ref{fig:chaining-runl-random-occurrences}
and~\ref{fig:chaining-runl-random-sizes} report only the initial
dynamic-programming call used by the special first-path loop.

\subsection{Does the bounded SIMD \texorpdfstring{$U$}{U} limit help?}
\label{sec:chaining-u-batch-experiment}
\begingroup
\small
\setlength{\intextsep}{5pt plus 1pt minus 1pt}

A standalone minimap2 v2.31 was instrumented with scalar shadows of bounded
SIMD B1 for \(L\in\{2,4,8,16,32,64\}\).  Each width recursively used its own
scores and the saved-anchor check from Algorithm~\ref{alg:chaining-b5}; the
shadows did not change minimap2's mapping.
For width \(L\), let \(E_L\) be the number of executed batches, \(Z_L\) the
executed batches for which no lane satisfies \(C_\lambda>H_\lambda\), and
\(Q_L\) the batches in the complete window without the \(U\) stop.  The two
pooled rates are
\[
 \operatorname{NoUpdate}_L=Z_L/E_L,
 \qquad
 \operatorname{Reduction}^{U}_L=(Q_L-E_L)/Q_L.
\]
The first rate counts outcomes that increment \(u\), while the second isolates
batch-loop iterations saved by \(U\).  Pooling sums counters before division,
so it is neither a median nor an unweighted mean.  The scalar-loop analogue of
the earlier \(\operatorname{Reduction}_L\) was 96.3--99.5\% for random
occurrences and 98.7--99.7\% for mosaic, but includes ordinary batching.
\(\operatorname{Reduction}^{U}_L\) isolates the counter's effect.

\begin{figure}[H]
\centering
% Generated by 06_chaining_u_batch_plot.py; do not hand-edit.
\begin{tikzpicture}[font=\scriptsize]
  \draw[draw=chainBlue!95!black,line width=.85pt] (2.10,3.75)--(2.52,3.75) node[anchor=west,text=black] {$L=2$};
  \draw[draw=chainBlue!95!black,line width=.85pt,densely dashed] (3.95,3.75)--(4.37,3.75) node[anchor=west,text=black] {$L=4$};
  \draw[draw=chainOrange!95!black,line width=.85pt] (5.80,3.75)--(6.22,3.75) node[anchor=west,text=black] {$L=8$};
  \draw[draw=chainOrange!95!black,line width=.85pt,densely dashed] (7.65,3.75)--(8.07,3.75) node[anchor=west,text=black] {$L=16$};
  \draw[draw=chainGreen!80!black,line width=.85pt] (9.50,3.75)--(9.92,3.75) node[anchor=west,text=black] {$L=32$};
  \draw[draw=chainGreen!80!black,line width=.85pt,densely dashed] (11.35,3.75)--(11.77,3.75) node[anchor=west,text=black] {$L=64$};
  \begin{scope}[shift={(1.15,.55)},x=.25cm,y=.027cm]
    \foreach \y in {0,25,50,75,100} {
      \draw[black!14] (0,\y)--(20,\y);
      \draw (-.18,\y)--(0,\y);
      \node[anchor=east,xshift=-1pt] at (0,\y) {\y};
    }
    \foreach \x in {0,5,10,15,20} {
      \draw (\x,0)--(\x,-3.3);
      \node[anchor=north,yshift=-1pt] at (\x,0) {\x};
    }
    \draw[->] (0,0)--(20.6,0);
    \draw[->] (0,0)--(0,105);
    \draw[draw=chainBlue!95!black,line width=.85pt] plot coordinates {(0.00,96.5162) (0.25,96.3613) (0.50,96.1963) (0.75,96.0490) (1.00,95.8422) (1.25,95.6649) (1.50,95.4567) (1.75,95.3126) (2.00,95.0661) (2.25,94.8689) (2.50,94.6063) (2.75,94.4151) (3.00,94.1489) (3.25,93.9304) (3.50,93.8001) (3.75,93.6301) (4.00,93.3730) (4.25,93.2353) (4.50,93.0552) (4.75,92.8565) (5.00,92.7629) (5.25,92.4635) (5.50,92.0786) (5.75,91.8551) (6.00,91.7345) (6.25,91.3997) (6.50,91.1061) (6.75,90.8319) (7.00,90.6916) (7.25,90.5679) (7.50,90.4128) (7.75,90.3142) (8.00,90.1198) (8.25,89.8699) (8.50,89.7075) (8.75,89.5332) (9.00,89.2000) (9.25,88.9205) (9.50,88.6366) (9.75,88.4840) (10.00,88.1310) (10.25,87.9766) (10.50,87.9381) (10.75,87.8470) (11.00,87.5873) (11.25,87.2918) (11.50,87.0654) (11.75,86.7934) (12.00,86.7934) (12.25,86.7544) (12.50,86.3533) (12.75,86.1008) (13.00,85.8657) (13.25,85.8657) (13.50,85.8516) (13.75,85.8079) (14.00,85.6683) (14.25,85.1285) (14.50,84.7932) (14.75,84.3618) (15.00,84.2897) (15.25,84.2897) (15.50,84.2261) (15.75,84.0232) (16.00,84.0232) (16.25,84.0963) (16.50,83.7346) (16.75,83.7346) (17.00,83.7346) (17.25,83.5290) (17.50,83.2431) (17.75,83.1580) (18.00,83.1580) (18.25,82.9968) (18.50,82.7799) (18.75,82.7030) (19.00,82.4893) (19.25,82.2480) (19.50,82.1706) (19.75,82.0027) (20.00,81.8374)};
    \draw[draw=chainBlue!95!black,line width=.85pt,densely dashed] plot coordinates {(0.00,96.0244) (0.25,95.8483) (0.50,95.6605) (0.75,95.4926) (1.00,95.2570) (1.25,95.0554) (1.50,94.8189) (1.75,94.6545) (2.00,94.3732) (2.25,94.1490) (2.50,93.8504) (2.75,93.6328) (3.00,93.3303) (3.25,93.0826) (3.50,92.9341) (3.75,92.7409) (4.00,92.4475) (4.25,92.2912) (4.50,92.0880) (4.75,91.8627) (5.00,91.7554) (5.25,91.4162) (5.50,90.9776) (5.75,90.7240) (6.00,90.5893) (6.25,90.2112) (6.50,89.8769) (6.75,89.5667) (7.00,89.4080) (7.25,89.2685) (7.50,89.0897) (7.75,88.9769) (8.00,88.7582) (8.25,88.4762) (8.50,88.2896) (8.75,88.0883) (9.00,87.7151) (9.25,87.4009) (9.50,87.0782) (9.75,86.9065) (10.00,86.5069) (10.25,86.3347) (10.50,86.2898) (10.75,86.1852) (11.00,85.8890) (11.25,85.5554) (11.50,85.3014) (11.75,84.9906) (12.00,84.9906) (12.25,84.9448) (12.50,84.4910) (12.75,84.2082) (13.00,83.9468) (13.25,83.9468) (13.50,83.9321) (13.75,83.8813) (14.00,83.7294) (14.25,83.1157) (14.50,82.7417) (14.75,82.2580) (15.00,82.1740) (15.25,82.1740) (15.50,82.1052) (15.75,81.8747) (16.00,81.8747) (16.25,81.9575) (16.50,81.5470) (16.75,81.5470) (17.00,81.5470) (17.25,81.3132) (17.50,80.9933) (17.75,80.8980) (18.00,80.8980) (18.25,80.7145) (18.50,80.4698) (18.75,80.3905) (19.00,80.1555) (19.25,79.8769) (19.50,79.7917) (19.75,79.6043) (20.00,79.4225)};
    \draw[draw=chainOrange!95!black,line width=.85pt] plot coordinates {(0.00,95.0438) (0.25,94.8258) (0.50,94.5947) (0.75,94.3855) (1.00,94.0934) (1.25,93.8431) (1.50,93.5492) (1.75,93.3471) (2.00,92.9974) (2.25,92.7185) (2.50,92.3488) (2.75,92.0804) (3.00,91.6989) (3.25,91.3970) (3.50,91.2090) (3.75,90.9701) (4.00,90.6129) (4.25,90.4210) (4.50,90.1745) (4.75,89.8875) (5.00,89.7598) (5.25,89.3320) (5.50,88.7914) (5.75,88.4864) (6.00,88.3229) (6.25,87.8549) (6.50,87.4418) (6.75,87.0627) (7.00,86.8603) (7.25,86.6972) (7.50,86.4794) (7.75,86.3356) (8.00,86.0598) (8.25,85.7218) (8.50,85.4933) (8.75,85.2333) (9.00,84.7834) (9.25,84.3886) (9.50,83.9956) (9.75,83.7913) (10.00,83.3185) (10.25,83.1119) (10.50,83.0410) (10.75,82.8997) (11.00,82.5361) (11.25,82.1076) (11.50,81.7927) (11.75,81.4252) (12.00,81.4252) (12.25,81.3857) (12.50,80.8370) (12.75,80.4997) (13.00,80.1841) (13.25,80.1841) (13.50,80.1691) (13.75,80.0914) (14.00,79.8671) (14.25,79.1464) (14.50,78.6982) (14.75,78.1028) (15.00,77.9812) (15.25,77.9812) (15.50,77.9064) (15.75,77.6241) (16.00,77.6241) (16.25,77.7198) (16.50,77.2199) (16.75,77.2199) (17.00,77.2199) (17.25,76.9614) (17.50,76.5946) (17.75,76.4729) (18.00,76.4729) (18.25,76.2581) (18.50,75.9288) (18.75,75.7609) (19.00,75.4994) (19.25,75.1194) (19.50,75.0608) (19.75,74.8382) (20.00,74.7254)};
    \draw[draw=chainOrange!95!black,line width=.85pt,densely dashed] plot coordinates {(0.00,94.0862) (0.25,93.8265) (0.50,93.5618) (0.75,93.3152) (1.00,92.9684) (1.25,92.6732) (1.50,92.3184) (1.75,92.0786) (2.00,91.6688) (2.25,91.3374) (2.50,90.8977) (2.75,90.5918) (3.00,90.1280) (3.25,89.7644) (3.50,89.5753) (3.75,89.2855) (4.00,88.8455) (4.25,88.6462) (4.50,88.3356) (4.75,88.0086) (5.00,87.8366) (5.25,87.3217) (5.50,86.7393) (5.75,86.3637) (6.00,86.1297) (6.25,85.6270) (6.50,85.1863) (6.75,84.7213) (7.00,84.4813) (7.25,84.2863) (7.50,83.9984) (7.75,83.8539) (8.00,83.5020) (8.25,83.1171) (8.50,82.9390) (8.75,82.6175) (9.00,82.0416) (9.25,81.6454) (9.50,81.2019) (9.75,80.9324) (10.00,80.4109) (10.25,80.1826) (10.50,80.0683) (10.75,79.8377) (11.00,79.3722) (11.25,78.9380) (11.50,78.5566) (11.75,78.1904) (12.00,78.1904) (12.25,78.1494) (12.50,77.4443) (12.75,77.1638) (13.00,76.8381) (13.25,76.8381) (13.50,76.7354) (13.75,76.6585) (14.00,76.2935) (14.25,75.5833) (14.50,74.9638) (14.75,74.3086) (15.00,74.1382) (15.25,74.1382) (15.50,74.0117) (15.75,73.7460) (16.00,73.7460) (16.25,73.8197) (16.50,73.3376) (16.75,73.3376) (17.00,73.3376) (17.25,73.1566) (17.50,72.5930) (17.75,72.4793) (18.00,72.4793) (18.25,72.1873) (18.50,71.8066) (18.75,71.5798) (19.00,71.2821) (19.25,70.7947) (19.50,70.7847) (19.75,70.6956) (20.00,70.5481)};
    \draw[draw=chainGreen!80!black,line width=.85pt] plot coordinates {(0.00,92.2100) (0.25,91.8641) (0.50,91.5134) (0.75,91.2100) (1.00,90.7383) (1.25,90.3608) (1.50,89.8887) (1.75,89.5751) (2.00,89.0744) (2.25,88.6181) (2.50,88.0772) (2.75,87.6511) (3.00,87.0448) (3.25,86.6168) (3.50,86.3446) (3.75,86.0258) (4.00,85.4203) (4.25,85.2036) (4.50,84.7513) (4.75,84.3697) (5.00,84.1707) (5.25,83.4918) (5.50,82.7024) (5.75,82.3885) (6.00,81.9623) (6.25,81.3435) (6.50,80.8304) (6.75,80.2239) (7.00,79.8652) (7.25,79.6881) (7.50,79.2237) (7.75,79.2544) (8.00,78.7304) (8.25,78.1956) (8.50,77.9972) (8.75,77.5428) (9.00,76.8402) (9.25,76.3099) (9.50,75.8742) (9.75,75.5749) (10.00,74.8987) (10.25,74.7745) (10.50,74.4922) (10.75,74.2745) (11.00,73.4463) (11.25,73.1188) (11.50,72.3823) (11.75,71.7339) (12.00,71.7339) (12.25,71.6201) (12.50,70.7196) (12.75,70.7221) (13.00,70.7110) (13.25,70.7110) (13.50,70.6121) (13.75,70.6089) (14.00,70.1781) (14.25,68.5479) (14.50,68.1955) (14.75,67.2721) (15.00,66.8632) (15.25,66.8632) (15.50,66.8641) (15.75,66.3891) (16.00,66.3891) (16.25,66.3903) (16.50,66.3348) (16.75,66.3348) (17.00,66.3348) (17.25,66.0780) (17.50,65.3531) (17.75,65.3469) (18.00,65.3469) (18.25,65.2794) (18.50,64.4307) (18.75,63.8822) (19.00,63.8088) (19.25,62.9626) (19.50,62.9756) (19.75,62.7682) (20.00,62.7934)};
    \draw[draw=chainGreen!80!black,line width=.85pt,densely dashed] plot coordinates {(0.00,84.7779) (0.25,83.9656) (0.50,83.3603) (0.75,82.8097) (1.00,81.8388) (1.25,81.1380) (1.50,80.2318) (1.75,79.8069) (2.00,78.5887) (2.25,77.7576) (2.50,76.6652) (2.75,76.1298) (3.00,75.0569) (3.25,74.3306) (3.50,73.3688) (3.75,72.8768) (4.00,71.9612) (4.25,71.4826) (4.50,70.8220) (4.75,70.2676) (5.00,69.3908) (5.25,68.7430) (5.50,66.8639) (5.75,66.4950) (6.00,64.8072) (6.25,63.6150) (6.50,62.9162) (6.75,62.9518) (7.00,62.4210) (7.25,61.8901) (7.50,60.6113) (7.75,60.6977) (8.00,59.6651) (8.25,58.3090) (8.50,57.7693) (8.75,56.9287) (9.00,56.6051) (9.25,55.6874) (9.50,55.4338) (9.75,54.4272) (10.00,52.7033) (10.25,52.4091) (10.50,52.4165) (10.75,51.9841) (11.00,51.0079) (11.25,50.3253) (11.50,48.9129) (11.75,47.9396) (12.00,47.9396) (12.25,47.5407) (12.50,44.3721) (12.75,44.3707) (13.00,44.3340) (13.25,44.3340) (13.50,43.9758) (13.75,44.0102) (14.00,43.2812) (14.25,39.5484) (14.50,39.5141) (14.75,36.9393) (15.00,35.4087) (15.25,35.4087) (15.50,35.4172) (15.75,34.3528) (16.00,34.3528) (16.25,34.3485) (16.50,34.1816) (16.75,34.1816) (17.00,34.1816) (17.25,34.1927) (17.50,34.2080) (17.75,34.2119) (18.00,34.2119) (18.25,34.2198) (18.50,34.2319) (18.75,34.2361) (19.00,34.2486) (19.25,34.2616) (19.50,34.2660) (19.75,34.2749) (20.00,34.2840)};
  \end{scope}
  \node[anchor=north west,font=\bfseries,fill=white,inner sep=1pt] at (1.23,3.20) {(a)};
  \node[anchor=north] at (3.65,.14) {divergence (\%)};
  \begin{scope}[shift={(7.85,.55)},x=.25cm,y=.027cm]
    \foreach \y in {0,25,50,75,100} {
      \draw[black!14] (0,\y)--(20,\y);
      \draw (-.18,\y)--(0,\y);
      \node[anchor=east,xshift=-1pt] at (0,\y) {\y};
    }
    \foreach \x in {0,5,10,15,20} {
      \draw (\x,0)--(\x,-3.3);
      \node[anchor=north,yshift=-1pt] at (\x,0) {\x};
    }
    \draw[->] (0,0)--(20.6,0);
    \draw[->] (0,0)--(0,105);
    \draw[draw=chainBlue!95!black,line width=.85pt] plot coordinates {(0.00,98.4268) (0.25,98.2985) (0.50,98.1908) (0.75,98.0890) (1.00,98.0006) (1.25,97.9251) (1.50,97.8505) (1.75,97.7892) (2.00,97.7314) (2.25,97.6835) (2.50,97.6428) (2.75,97.6029) (3.00,97.5639) (3.25,97.5321) (3.50,97.4964) (3.75,97.4683) (4.00,97.4419) (4.25,97.4200) (4.50,97.3949) (4.75,97.3696) (5.00,97.3517) (5.25,97.3303) (5.50,97.3107) (5.75,97.2918) (6.00,97.2751) (6.25,97.2599) (6.50,97.2457) (6.75,97.2313) (7.00,97.2175) (7.25,97.2049) (7.50,97.1904) (7.75,97.1841) (8.00,97.1687) (8.25,97.1558) (8.50,97.1430) (8.75,97.1312) (9.00,97.1198) (9.25,97.1098) (9.50,97.0991) (9.75,97.0897) (10.00,97.0784) (10.25,97.0771) (10.50,97.0670) (10.75,97.0612) (11.00,97.0556) (11.25,97.0495) (11.50,97.0423) (11.75,97.0389) (12.00,97.0353) (12.25,97.0264) (12.50,97.0221) (12.75,97.0201) (13.00,97.0124) (13.25,97.0103) (13.50,97.0058) (13.75,97.0024) (14.00,96.9988) (14.25,96.9948) (14.50,96.9887) (14.75,96.9872) (15.00,96.9868) (15.25,96.9839) (15.50,96.9814) (15.75,96.9781) (16.00,96.9785) (16.25,96.9754) (16.50,96.9723) (16.75,96.9669) (17.00,96.9629) (17.25,96.9603) (17.50,96.9608) (17.75,96.9569) (18.00,96.9547) (18.25,96.9533) (18.50,96.9539) (18.75,96.9512) (19.00,96.9500) (19.25,96.9503) (19.50,96.9470) (19.75,96.9487) (20.00,96.9484)};
    \draw[draw=chainBlue!95!black,line width=.85pt,densely dashed] plot coordinates {(0.00,98.2029) (0.25,98.0303) (0.50,97.8830) (0.75,97.7407) (1.00,97.6138) (1.25,97.5072) (1.50,97.4027) (1.75,97.3145) (2.00,97.2331) (2.25,97.1638) (2.50,97.1096) (2.75,97.0545) (3.00,97.0045) (3.25,96.9601) (3.50,96.9117) (3.75,96.8755) (4.00,96.8405) (4.25,96.8121) (4.50,96.7794) (4.75,96.7501) (5.00,96.7292) (5.25,96.7074) (5.50,96.6841) (5.75,96.6625) (6.00,96.6442) (6.25,96.6267) (6.50,96.6151) (6.75,96.5977) (7.00,96.5882) (7.25,96.5771) (7.50,96.5596) (7.75,96.5485) (8.00,96.5320) (8.25,96.5208) (8.50,96.5066) (8.75,96.5018) (9.00,96.4962) (9.25,96.4902) (9.50,96.4801) (9.75,96.4769) (10.00,96.4680) (10.25,96.4719) (10.50,96.4652) (10.75,96.4604) (11.00,96.4651) (11.25,96.4655) (11.50,96.4631) (11.75,96.4630) (12.00,96.4643) (12.25,96.4624) (12.50,96.4632) (12.75,96.4678) (13.00,96.4642) (13.25,96.4669) (13.50,96.4705) (13.75,96.4696) (14.00,96.4702) (14.25,96.4693) (14.50,96.4662) (14.75,96.4734) (15.00,96.4775) (15.25,96.4785) (15.50,96.4784) (15.75,96.4785) (16.00,96.4834) (16.25,96.4846) (16.50,96.4857) (16.75,96.4830) (17.00,96.4813) (17.25,96.4808) (17.50,96.4829) (17.75,96.4811) (18.00,96.4809) (18.25,96.4803) (18.50,96.4833) (18.75,96.4832) (19.00,96.4831) (19.25,96.4854) (19.50,96.4824) (19.75,96.4860) (20.00,96.4870)};
    \draw[draw=chainOrange!95!black,line width=.85pt] plot coordinates {(0.00,97.7558) (0.25,97.4803) (0.50,97.2311) (0.75,96.9938) (1.00,96.7796) (1.25,96.6016) (1.50,96.4325) (1.75,96.2825) (2.00,96.1431) (2.25,96.0295) (2.50,95.9356) (2.75,95.8388) (3.00,95.7546) (3.25,95.6826) (3.50,95.6118) (3.75,95.5574) (4.00,95.4976) (4.25,95.4524) (4.50,95.4057) (4.75,95.3581) (5.00,95.3170) (5.25,95.2886) (5.50,95.2528) (5.75,95.2197) (6.00,95.1876) (6.25,95.1559) (6.50,95.1490) (6.75,95.1286) (7.00,95.1225) (7.25,95.1176) (7.50,95.0941) (7.75,95.0988) (8.00,95.0846) (8.25,95.0622) (8.50,95.0431) (8.75,95.0477) (9.00,95.0544) (9.25,95.0684) (9.50,95.0738) (9.75,95.0857) (10.00,95.0801) (10.25,95.0996) (10.50,95.1024) (10.75,95.1019) (11.00,95.1195) (11.25,95.1344) (11.50,95.1550) (11.75,95.1699) (12.00,95.1895) (12.25,95.2065) (12.50,95.2235) (12.75,95.2435) (13.00,95.2486) (13.25,95.2629) (13.50,95.2760) (13.75,95.2883) (14.00,95.3021) (14.25,95.3143) (14.50,95.3154) (14.75,95.3375) (15.00,95.3553) (15.25,95.3666) (15.50,95.3757) (15.75,95.3794) (16.00,95.4064) (16.25,95.4170) (16.50,95.4302) (16.75,95.4337) (17.00,95.4397) (17.25,95.4529) (17.50,95.4622) (17.75,95.4653) (18.00,95.4743) (18.25,95.4797) (18.50,95.4898) (18.75,95.4936) (19.00,95.4980) (19.25,95.5054) (19.50,95.5050) (19.75,95.5134) (20.00,95.5196)};
    \draw[draw=chainOrange!95!black,line width=.85pt,densely dashed] plot coordinates {(0.00,97.3124) (0.25,96.8958) (0.50,96.5090) (0.75,96.1400) (1.00,95.7867) (1.25,95.4924) (1.50,95.2029) (1.75,94.9580) (2.00,94.7056) (2.25,94.5045) (2.50,94.2995) (2.75,94.1078) (3.00,93.9542) (3.25,93.8301) (3.50,93.6955) (3.75,93.5879) (4.00,93.4972) (4.25,93.3911) (4.50,93.3013) (4.75,93.2026) (5.00,93.1451) (5.25,93.0942) (5.50,93.0291) (5.75,92.9710) (6.00,92.9110) (6.25,92.8700) (6.50,92.8723) (6.75,92.8410) (7.00,92.8363) (7.25,92.8022) (7.50,92.7944) (7.75,92.7742) (8.00,92.7847) (8.25,92.7434) (8.50,92.7582) (8.75,92.8139) (9.00,92.8242) (9.25,92.8751) (9.50,92.9259) (9.75,92.9200) (10.00,92.9293) (10.25,92.9968) (10.50,93.0153) (10.75,93.0148) (11.00,93.0874) (11.25,93.1329) (11.50,93.1817) (11.75,93.2114) (12.00,93.2895) (12.25,93.3397) (12.50,93.4010) (12.75,93.4520) (13.00,93.4892) (13.25,93.5345) (13.50,93.5768) (13.75,93.6123) (14.00,93.6551) (14.25,93.6939) (14.50,93.6969) (14.75,93.7532) (15.00,93.7895) (15.25,93.8089) (15.50,93.8384) (15.75,93.8612) (16.00,93.9241) (16.25,93.9606) (16.50,93.9884) (16.75,94.0155) (17.00,94.0459) (17.25,94.0837) (17.50,94.1060) (17.75,94.1289) (18.00,94.1518) (18.25,94.1708) (18.50,94.1982) (18.75,94.2186) (19.00,94.2352) (19.25,94.2617) (19.50,94.2687) (19.75,94.2883) (20.00,94.3056)};
    \draw[draw=chainGreen!80!black,line width=.85pt] plot coordinates {(0.00,96.4309) (0.25,95.6378) (0.50,94.9193) (0.75,94.2149) (1.00,93.5748) (1.25,93.0034) (1.50,92.4752) (1.75,91.9855) (2.00,91.5352) (2.25,91.2091) (2.50,90.8234) (2.75,90.4868) (3.00,90.1665) (3.25,89.9360) (3.50,89.6857) (3.75,89.4776) (4.00,89.3156) (4.25,89.1134) (4.50,88.9850) (4.75,88.8915) (5.00,88.7196) (5.25,88.5828) (5.50,88.4479) (5.75,88.2980) (6.00,88.2218) (6.25,88.2109) (6.50,88.1793) (6.75,88.1632) (7.00,88.1110) (7.25,88.0819) (7.50,88.0840) (7.75,87.9982) (8.00,88.1268) (8.25,88.0274) (8.50,88.0752) (8.75,88.1195) (9.00,88.2056) (9.25,88.3866) (9.50,88.4918) (9.75,88.5215) (10.00,88.5183) (10.25,88.6110) (10.50,88.6624) (10.75,88.7250) (11.00,88.8811) (11.25,89.0119) (11.50,89.1069) (11.75,89.1754) (12.00,89.3309) (12.25,89.4390) (12.50,89.4975) (12.75,89.6586) (13.00,89.6849) (13.25,89.7612) (13.50,89.8211) (13.75,89.8931) (14.00,89.9766) (14.25,90.0709) (14.50,90.1167) (14.75,90.2447) (15.00,90.2977) (15.25,90.3553) (15.50,90.4040) (15.75,90.4554) (16.00,90.6066) (16.25,90.6921) (16.50,90.7547) (16.75,90.8137) (17.00,90.8735) (17.25,90.9726) (17.50,91.0156) (17.75,91.0435) (18.00,91.1076) (18.25,91.1590) (18.50,91.2275) (18.75,91.2514) (19.00,91.2988) (19.25,91.3678) (19.50,91.3900) (19.75,91.4476) (20.00,91.4987)};
    \draw[draw=chainGreen!80!black,line width=.85pt,densely dashed] plot coordinates {(0.00,92.9185) (0.25,90.0130) (0.50,87.5779) (0.75,85.4987) (1.00,83.9055) (1.25,82.4383) (1.50,81.5247) (1.75,80.4936) (2.00,79.6053) (2.25,78.8733) (2.50,78.1753) (2.75,77.4098) (3.00,76.6591) (3.25,76.1064) (3.50,75.6038) (3.75,75.2965) (4.00,74.1380) (4.25,73.8884) (4.50,73.5978) (4.75,73.2714) (5.00,73.1984) (5.25,72.7856) (5.50,72.3575) (5.75,72.3579) (6.00,72.0691) (6.25,71.9457) (6.50,71.8567) (6.75,71.9030) (7.00,71.6708) (7.25,71.5771) (7.50,71.6597) (7.75,71.4496) (8.00,71.6689) (8.25,71.3064) (8.50,71.5220) (8.75,71.7150) (9.00,71.7225) (9.25,71.7354) (9.50,72.1242) (9.75,72.2876) (10.00,72.2667) (10.25,72.2718) (10.50,72.3351) (10.75,72.8950) (11.00,73.1483) (11.25,73.2508) (11.50,73.3641) (11.75,73.7356) (12.00,74.1856) (12.25,74.3467) (12.50,74.7109) (12.75,75.2051) (13.00,75.1270) (13.25,75.4529) (13.50,75.7771) (13.75,75.9520) (14.00,76.3893) (14.25,76.8061) (14.50,77.0921) (14.75,77.3720) (15.00,77.5040) (15.25,77.9081) (15.50,77.9447) (15.75,78.2137) (16.00,78.5387) (16.25,78.5548) (16.50,78.8209) (16.75,79.0279) (17.00,79.4126) (17.25,79.7491) (17.50,79.9021) (17.75,79.9410) (18.00,80.1245) (18.25,80.3760) (18.50,80.4480) (18.75,80.5721) (19.00,80.7306) (19.25,80.9572) (19.50,81.0372) (19.75,81.1744) (20.00,81.3284)};
  \end{scope}
  \node[anchor=north west,font=\bfseries,fill=white,inner sep=1pt] at (7.93,3.20) {(b)};
  \node[anchor=north] at (10.35,.14) {divergence (\%)};
  \node[rotate=90,anchor=south] at (.10,1.90) {$\operatorname{Reduction}^{U}_{L}$ (\%)};
\end{tikzpicture}

\caption{Pooled \(U\)-only reduction at each divergence for (a) random inputs
(six seeds) and (b) mosaic inputs (three seeds).  Each point is a ratio of
summed integer counters, rather than a mean of per-input percentages.}
\label{fig:chaining-u-batch-reduction}
\end{figure}

\begin{table}[H]
\centering
\caption{\(\operatorname{NoUpdate}_L/\operatorname{Reduction}^{U}_L\) (percent).
The 0\% and 10\% rows pool seeds at that level.  Pooled denotes the full sweep.
only Random and Mosaic vary divergence.  Random uses \(N=4000\), a 4.5-Mbase
target, 997--1002-base queries, and 81 substitution levels, with ten exact
30-base intervals.  Mosaic uses \(T=Q=U_0AU_1A'_\delta U_2\), length 67072.
Spaced copies use \(\lvert S\rvert=4096\), \(G=16384\), and
\(R\in\{1,2,4,8,16,32,64\}\).  Coordinate control uses
\(m\in\{1,2,4,\ldots,1024\}\).  Defaults are \(k=15\), \(w=10\), \(S=25\), and
\texttt{max\_chain\_iter}=5000.}
\label{tab:chaining-u-batch-rates}
\scriptsize
\setlength{\tabcolsep}{2.5pt}
\renewcommand{\arraystretch}{1.0}
\resizebox{\linewidth}{!}{%
\begin{tabular}{@{}llrrrrrr@{}}
\toprule
\textbf{Subset} & \textbf{Input} & \(L=2\) & \(L=4\) & \(L=8\) & \(L=16\) & \(L=32\) & \(L=64\)\\
\midrule
% Generated by 06_chaining_u_batch_plot.py; do not hand-edit.
0\% divergence & Random occurrences & 92.9/96.5 & 87.6/96.0 & 80.1/95.0 & 66.8/94.1 & 50.1/92.2 & 50.1/84.8\\
 & Mosaic & 92.9/98.4 & 87.5/98.2 & 80.0/97.8 & 66.6/97.3 & 49.9/96.4 & 49.7/92.9\\
\cmidrule(lr){1-8}
10\% divergence & Random occurrences & 93.0/88.1 & 87.7/86.5 & 80.3/83.3 & 67.1/80.4 & 50.5/74.9 & 50.5/52.7\\
 & Mosaic & 91.6/97.1 & 82.8/96.5 & 68.9/95.1 & 47.3/92.9 & 26.1/88.5 & 21.2/72.3\\
\cmidrule(lr){1-8}
Pooled & Random occurrences & 93.0/92.9 & 87.7/92.0 & 80.2/90.0 & 67.0/88.2 & 50.4/84.6 & 50.4/70.2\\
 & Mosaic & 91.7/97.4 & 83.1/96.9 & 69.7/95.7 & 48.8/94.0 & 27.7/90.5 & 23.0/77.4\\
 & Spaced copies & 97.3/99.3 & 92.9/99.1 & 82.2/98.7 & 62.6/98.5 & 35.4/97.9 & 15.4/90.6\\
 & Coordinate control & 93.2/53.1 & 88.0/47.4 & 80.7/36.8 & 67.6/28.8 & 51.0/17.8 & 27.0/0.0

\\\bottomrule
\end{tabular}}
\end{table}

All 761 queries had validated PAF matches.  For increasing widths, \(U\) is
\(13,7,4,2,1,1\).  Omitted batches per random-arm check were 13.17, 7.44, and 2.36 at
\(L=2,16,64\); at 20\% divergence the reductions were 81.8, 70.5, and 34.3\%.
Mosaic and spaced copies saved more.  Coordinate control at \(L=64\) was the
only zero.  The requested decreasing-score/decreasing-position tie order was
audited, but end sorting follows forward scoring and cannot affect these counts.
The data therefore support retaining \(U\).  This measures opportunity rather
than wall-clock speedup.  Final cost depends on the SIMD reduction for
\(\operatorname{Any}(\mathbf I)\).
\endgroup


\section{BioseqDB \texttt{max\_drop} analysis}
\label{sec:chaining-max-drop-analysis}

The BioseqDB backtracking procedure follows predecessor links until it reaches
the sentinel \(p_j=-1\) or an anchor that was already assigned to an earlier
chain.  Minimap2 can stop a predecessor scan earlier when its running score
falls sufficiently far below the best score seen on that scan.  This
\texttt{max\_drop} rule introduces score arithmetic, comparisons, and
running-maximum and lane-stop state during scan steps.  It can therefore help
only when the steps that it avoids repay the cost of those checks.  This
section measures that opportunity on minimap2 chaining data.  It does not
measure BioseqDB running time.

In the two synthetic repeated-copy screens below, the default \(D=500\)
cutoff did not avoid a single backtracking scan step.  The separate four-lane
\(\operatorname{Run}_4\) condition was nevertheless frequent, because it tests
path adjacency rather than score drawdown.

The same experiment also measures the condition used by the special
BioseqDB first-path branch, \(\operatorname{Run}_L\), where \(L\) is the SIMD
batch width.  The two questions are kept separate.  A score drop can shorten
a predecessor scan.  In contrast, \(\operatorname{Run}_L\) processes a
consecutive part of the highest-ranked predecessor path in a vector batch.

\subsection{Measured quantities}

Let \(f_i\) be minimap2's dynamic-programming score for the best chain ending
at anchor \(i\).  The score plots include every anchor in the initial-DP array
of each retained observation.  Candidate ends are the subset satisfying
\(f_i\geq\texttt{min\_sc}\), where the default value in these initial chaining
calls was 40.  Minimap2 normally orders these ends by score.  The experiment
ordered them by increasing \((f_i,i)\) and consumed the array backwards.  It
therefore used decreasing score and, on a score tie, decreasing anchor-array
index.  Here ``position'' means the index in minimap2's coordinate-sorted
anchor array.  The anchor array itself was not reordered.

One executed search is a predecessor scan started from an eligible chain end
whose anchor is unmarked on entry.  For end \(e\), let \(B\) be the number of
steps made by the BioseqDB-style baseline scan, which stops only at
\(p_j=-1\) or at a previously marked anchor.  Let \(b_m\) be the boundary
reached after predecessor step \(m\), and define
\[
  s_m=
  \begin{cases}
    f_e-f_{b_m}, & b_m\geq0,\\
    f_e, & b_m=-1,
  \end{cases}
  \qquad
  M_m=\max(0,s_1,\ldots,s_m).
\]
Define the score drop at step \(m\) as
\[
  d_m=M_m-s_m,
  \qquad
  \Delta(e)=\max_{1\leq m\leq B}d_m.
\]
Thus \(d_m\) is the distance from the best score already seen in this
backtrack to the score at the current step.  The value \(\Delta(e)\) is the
largest such drop in the scan from end \(e\).  The max-drop predicate with
threshold \(D\) fires exactly when \(d_m>D\).  Equality does not fire.

Let \(I(D)\) be the steps made from the same entry
marking state when the literal minimap2 rule stops at the first predicate
firing.  If the predicate never fires, then \(I(D)=B\).  The local suffix and
the per-search local suffix fraction are
\begin{equation}
  S(D)=B-I(D),
  \qquad
  \beta(D)=
  \begin{cases}
    1-I(D)/B, & B>0,\\
    0, & B=0.
  \end{cases}
\label{eq:chaining-max-drop-benefit}
\end{equation}
The fixed-start calculation uses the same chain end and the same incoming
marks as the baseline scan, and it does not write new marks.  Thus \(S(D)\)
is the part of that baseline path after its first crossing.  Across several
scans, the reported fraction is \(\sum S(D)/\sum B\).  A productive crossing is one for which
\(S(D)>0\), not merely one on which the predicate becomes true.  This
distinction matters because the predicate can fire on the final scan step and
leave no suffix.

This one-scan calculation need not equal the work of applying the rule to a
complete backtrack.  Committing a shorter path changes the marked anchors and
can expose later candidate ends.  Let \(\widetilde E\) and \(\widetilde I\)
denote the searches and backtracking scan steps when every cutoff is applied
and the marks are updated.  The experiment reports both calculations.  Steps
after a crossing are not called saved work unless the complete replay also
keeps search ownership and accepted output unchanged.

The complete replay retains minimap2's best-cut selection after a path, even
when no score-drop cutoff fires.  A matched no-drop replay keeps that
bookkeeping but disables only the score drop.  Comparing the finite and
no-drop replays therefore isolates the
\(D\)-threshold decision.  Comparing either replay with the BioseqDB-style
baseline can additionally expose changes caused by the independent best-cut
boundary.  The audit checks search counts, scan-step counts, accepted chains,
execution traces, and final marks separately.  Here accepted-chain output is
the output of the initial-DP backtracking replay, not final PAF output.

The cost model makes the break-even point explicit.  Let \(v>0\) be the cost
of one backtracking scan step without the new rule, let \(h\geq0\) be the
incremental cost of one drop check, and let \(C\) be the number of checks.
For a conservative implementation budget, one logical score-check exposure
is charged per limited scan step, so \(C=\sum I\).  This is not a measured
machine-instruction count, nor the literal count of native threshold
comparisons.  A scan step that establishes a new maximum may bypass that
comparison.  Ignoring fixed setup cost, the fixed-start model is
\begin{equation}
  T_0=v\sum B,
  \qquad
  T_D^{\mathrm{fixed}}=v\sum I+hC.
\label{eq:chaining-max-drop-cost}
\end{equation}
Writing \(r=h/v\), and noting that \(C>0\) for every reported exposure, the
local opportunity can repay the checks only if
\begin{equation}
  r<\frac{\sum S}{C}.
\label{eq:chaining-max-drop-break-even}
\end{equation}
Equivalently, it requires
\(\sum S/\sum B>r/(1+r)\).  The threshold \(D\) controls where a scan stops.
It is not the instruction-cost ratio \(r\).  For the complete replay,
however,
\begin{equation}
  T_D^{\mathrm{full}}=v\widetilde I+h\widetilde C,
\label{eq:chaining-max-drop-full-cost}
\end{equation}
where \(\widetilde C\) is its own check exposure.  Consequently, even a large
fixed-start suffix is insufficient when \(\widetilde I\geq\sum B\), or when the
cutoff changes an output that BioseqDB must preserve.
This expression models variable scan work and excludes per-search setup cost.

For the SIMD replay, later candidate ends are collected in batches containing
at most \(L\) candidates, with \(L\in\{4,8,16\}\).  Each batch uses one common
pre-batch marking state.  The baseline
and drop-limited depth of each lane are recorded.  The batch critical depth
is the maximum lane depth, so a productive lane helps the critical path only
when it shortens that maximum.  These are logical work counters, not SIMD
instructions or clock cycles.

For the highest-ranked first path, let \(j\) be the current anchor, \(c\) the
number of anchors already consumed, and \(n\) the anchor-array length.  The
measured BioseqDB predicate is
\begin{equation}
  \operatorname{Run}_L(j,c)
  \iff
  j\geq L-1
  \land c\leq n-L
  \land
  \bigwedge_{u=0}^{L-1}p_{j-u}=j-u-1.
\label{eq:chaining-max-drop-runl}
\end{equation}
The relation \(p_0=-1=0-1\) may complete the final successful batch.  If
\(B_L\) is the number of true evaluations, \(E_L\) is the number of all loop
evaluations, and \(q\) is the scalar path length, the two reported quantities
are
\begin{equation}
  \operatorname{Fire}_L=\frac{B_L}{E_L},
  \qquad
  \operatorname{Reduction}_L
  =\frac{q-E_L}{q}
  =\frac{(L-1)B_L}{q}.
\label{eq:chaining-max-drop-runl-metrics}
\end{equation}
The symbol \(B_L\) counts successful fast batches and is unrelated to the
baseline scan-step count \(B\) above.  The reported data have \(E_L>0\) and
\(q>0\).  The second value is the ideal fraction of scalar loop evaluations
removed by the mixed fast/scalar loop.  It is not a measured speedup.

\subsection{Experimental design and validation}

The experiment used a disposable copy of minimap2 2.31-r1302, upstream tag
\texttt{v2.31} at commit \texttt{3c28777}.  It was built and run independently
of the BioseqDB build system.  Mapping used the default \texttt{map-ont} preset
and one mapping thread.  The experimental changes only emitted telemetry and
applied the score/index tie order described above.  When telemetry was
disabled, the modified executable and the official upstream executable
produced byte-identical PAF output in four fixed equivalence checks.

An experimental arm fixes the reference geometry.  The primary arm used a
4.5-Mbp reference.  An observation is one mapped query for a fixed arm, input
block, \(N\), and divergence.  A cell fixes the arm, \(N\), and divergence and
contains six observations when all are feasible.  The experiment used six
independently seeded input blocks with query lengths 997, 998, 999, 1000,
1001, and 1002 bases.  One query-length value is associated with one generator
seed, so this design does not separate query-length effects from seed
variation.  References across \(N\) and queries across divergence are nested
within each block.  The planted homologous-locus count was
\[
  N\in\{1,2,3,4,5,10,20,30,40,50,100,200,300,400,500,
       1000,2000,3000,4000\}.
\]
For each \(N\), the requested query substitution percentage ranged from 0 to 10 percent in
0.25-percentage-point steps.  Divergence used nested substitution-only
prefixes, and the copies used nested placement prefixes within each generated
input.  The 4,674 primary observations are therefore repeated measurements
in six input blocks, not 4,674 independent replicates.  Ratios were calculated
within an input before the median, linearly interpolated interquartile range,
and observed range were taken across the six inputs.  These ranges are
descriptive, not confidence intervals.

Five thousand non-overlapping copies of even the shortest query require
4,985,000 bases, before any separating sequence is added.  The requested
\(N=5000\) condition is therefore structurally impossible in an exactly
4.5-Mbp reference.  The canonical data record its 246 planned primary
observations as structural missing values rather than as zeros.  A separate
5.625-Mbp sensitivity arm contains \(N=5000\).  It changes reference length
together with \(N\), so it is not an extension of the primary \(N\)-axis.

Before the full sweep, a corner pilot fixed the validity rule at one forward
PAF record that overlaps a planted locus and covers at least 50 percent of
both the query and the intended copy.  This replaced an initial 90-percent
rule because minimap2's anchor-only PAF spans can be shorter than a base-level
alignment even for a clear planted-locus match.  All 4,920 feasible
observations passed the final rule.  The lowest observed best overlap was
80.34 percent.  The stricter 90-percent descriptive check retained 4,560 of
4,920 observations.  Repeated copies make the mapping non-unique, so this
gate proves a planted-locus match, not unique-locus recovery.

Score ties occurred in 4,674 of 4,920 initial calls.  At least one equal-score
group was reordered in 4,594 calls, although the highest-ranked anchor did
not change in this workload.  Thus the requested tie rule was active for
lower-ranked candidates even though it did not change the first endpoint.
An independent full-data audit accepted all 4,920 feasible observations plus
246 structural records with no error or warning.  A second oracle
reconstructed 3,216,718 dumped anchors and 3,113,498 candidate ends in 14
fixed cases selected before result inspection.  It
matched the score histograms, scalar drop rule, all three
\(\operatorname{Run}_L\) widths, and the SIMD batch counters exactly.

The seed used in the corner pilot remains one of the six fixed inputs.  A
leave-pilot-seed-out sensitivity left five input blocks.  The default-drop
metrics remained identically zero over 300,886,632 checks, the \(N=4000\)
median score still fell from 501 through 294 to 153 at 0, 5, and 10 percent
divergence, and every retained input still contained a successful
\(\operatorname{Run}_L\) batch at each width.  The conclusions below therefore
do not depend on the viewed pilot input.

\subsection{DP-score distribution across anchors}

Figure~\ref{fig:chaining-max-drop-score-counts} shows the requested count of
initial-DP anchors at each exact integer score.  A point is the arithmetic
mean of the six input-level counts, including zero counts, so the points in a
panel sum to that panel's mean DP-anchor count.  No score-width binning or
smoothing is used.  Since changing \(N\) did not visibly change the
distribution shape, the figure shows only the larger \(N=4000\) case.  The
orange line marks \(\texttt{min\_sc}=40\).

\begin{figure}[H]
\centering
\begin{CacheMeCode}{python, custom include command={\input{\robExtAddCachePathAndName{\robExtFinalHash.pgf}}}, set placeholder eval={__LINEWIDTH__}{\lenToCmNoUnit[in]{\linewidth}}}
import matplotlib
matplotlib.use("pgf")
import matplotlib.pyplot as plt

matplotlib.rcParams.update({
    "font.family": "serif",
    "font.serif": [],
    "text.usetex": True,
})

def expand_ranges(spec):
    values = []
    for token in spec.split(","):
        if "-" in token:
            lo, hi = (int(value) for value in token.split("-"))
            values.extend(range(lo, hi + 1))
        else:
            values.append(int(token))
    return values

# Each tuple stores a combined six-input count followed by the exact scores
# having that count.  Dividing by six gives the plotted equal-input mean.
score_groups = {
    0: [
        (4000, "18,20-21,23,26,31,33,35-36,38,40,43,45,47-50,52,55,61-62,66,68,71,73,77-78,82-84,87-89,91-95,97-98,100,102,107,109-110,113,115,119,121,125,128,130,132,136-137,140,145,151,154-155,163,168-169,171,173-174,176,180,183-185,187,191,193,200-202,206-207,209,211,214-216,218,222,225,230,232,235,237-239,243-244,248-249,251,256,258,261-262,264,268,270,276-284,287,291-292,295,302-307,310,314,317-318,321,326,328-330,335-336,341,346,350,354,356-359,365-367,370,372,375,380-381,384,387,392-393,396,399-400,402,404,406-407,411,413-414,416,418,420,422,424-426,429,435,437,440,443-444,446-447,452,454,456,460,463,465,467-468,470,472-475,478,482,484,486,488,491,493,495,498-500,502-503,505-506,508-509,511,518,524,528-529,533,537-539,541,544,548,552-553,556,561,566,572,574,578-582,586-587,592-595,598,605,608,610,614,616-617,619,622,629,632,637-638,641,643-645,649,652,654-656,661-662,665,667-668,673,675,677,679-680,682,685,688,692,695,697,699-700,702,705-709,712-714,716,720,723,727-728,731,737-738,749,756,758-759,761,763-764,766,768-769,771-772,775,777,782-783,786,789,791,793,795-796,798,802,808,813,816-817,821,825-826,828,830,832-833,837,839,841,843,845,848,854,856-857,859,862,865,868-869,874,876,880,884-885,887,889-895,898,903-905,910,924-925,927-928,931-932,935-937,940-941,946,950-951,954-958,960,963,970,974,977,979-988,990-991,995-996"),
        (8000, "22,24-25,27,30,39,44,54,60,67,72,79-80,85,99,104,108,111,117-118,123,129,131,139,141-142,144,157-158,160-161,166,170,175,178,182,188-189,192,194-196,203-204,210,212-213,223,226,229,231,234,236,240,247,252-253,257,263,273-275,285,288,293-294,296,311,315,319,324,327,332,340,347,355,363-364,373-374,376-377,382-383,389,391,395,398,401,409-410,412,419,431,436,438-439,448-449,453,457-459,462,466,469,479-481,487,489,496,501,504,512,516,519,521,525,527,530,534-535,540,554,558,563,568,573,575,589,596,601,606,611,615,621,625,630-631,633,635,639,650-651,653,658-659,666,681,684,686,689,693,698,703,715,717,719,721,733,736,739,742,744,747-748,751,753-754,760,774,776,779-781,785,792,801,806-807,810-812,814,819,822-823,827,829,834,836,838,842,853,855,861,864,867,871,873,879,882-883,899,902,906-909,913,916,919-920,944,948,959,961,966,968,975-976,978"),
        (8001, "669"),
        (12000, "34,53,69,74,90,120,134,153,220,224,245-246,297,322-323,343-345,352,386,423,430,442,477,497,514,543,547,555,564,571,584,588,597,609,634,647,657,672,674,696,784,794,803,844,852,860,875,896,918,922,930,933"),
        (16000, "59,103,116,147,265,312,333,362,620,636,726,734,767,943,945,965"),
        (20000, "743"),
        (24001, "15"),
    ],
    5: [
        (4000, "20-21,24-25,28,31-33,36-38,40,42,46-47,50,58,62-63,65,68-70,74-76,86,93,95,97,100,106-107,109,112,118-119,124,129,138,140,142-143,145,151-152,155,158-159,164-165,167,169,173,180,183,185,191-192,195,197,200,202-203,206-207,209,214-215,220-222,224,228,235-237,239,242-244,249-250,254,261,269-272,278,280,283,285,292-293,295-297,300-302,304,306,309-310,313-316,318-319,321-322,324,328,333,337-338,342-343,351-352,354,356-358,361-362,364,366,368,375-376,381,383-385,389,391-392,394-395,402-403,409,411,415-416,418,420-422,427,430-431,433,435-436,442,445-448,452,455,465,467-468,472,475,477,479-480,482,484,491,493,498,500-501,503-504,510,513,515-516,518,520-523,528,530,532,534,536,543-545,547,550,552,559-560,565,568-570,572,579,590,596,600,603,606,611,613,618"),
        (8000, "43,45,53,56,60,73,77,88,92,104,111,120,122-123,126-127,131-133,137,144,148,157,162-163,168,170-172,174,186,204,210,218-219,223,229,234,240,251-252,264-265,274,282,312,317,325,329,332,334,340-341,345,350,355,365,367,370,378-379,387,393,399,407-408,413-414,425,429,440,450-451,462,466,470-471,478,485,488,490,497,499,514,525,529,551,563,578"),
        (12000, "30,48,61,85,103,146-147,189,259,279,290,294,327,349,353,371,483,489,527,542,557,585"),
        (16000, "255,299"),
        (24001, "15"),
    ],
    10: [
        (4000, "16-18,20-21,26-29,32,35-36,41,43-45,47,50-51,57-58,62,64-65,68-70,75-76,80-82,88-89,92-93,99-103,105,107-108,118,121,123,126,131,133,135,138,140,145,147-148,150,153,159,162-163,166-168,173,175,177,179-180,190,192,196,198-201,203-204,209,211-212,215-216,218-219,221-224,226-227,230,232,236,242-244,248,251,255,257,259,261,263-264,274,276-277,287,299,310,314,325,328-329,331,338,343-344,349,355"),
        (8000, "30,56,59,66,74,90,111,114,117,132,136,144,151,155,170,178,182,185,188,195,197,231,237,239,246,258,273,278,285,288,295"),
        (12000, "85,96,129,194"),
        (24000, "15"),
    ],
}
mean_totals = {0: 722667, 5: 317333.5, 10: 134000}

fig, axes = plt.subplots(1, 3, sharex=True, sharey=True)
fig.set_size_inches(0.94 * __LINEWIDTH__, 0.36 * __LINEWIDTH__, forward=True)
for ax, divergence in zip(axes, (0, 5, 10)):
    for combined_count, encoded_scores in score_groups[divergence]:
        scores = expand_ranges(encoded_scores)
        ax.scatter(scores, [combined_count / 6.0] * len(scores),
                   s=3.0, color="#4C78A8", linewidths=0)
    ax.axvline(40, color="#F28E2B", linestyle="--", linewidth=0.9)
    ax.set_xlim(0, 1000)
    ax.set_ylim(0, 4300)
    ax.set_xticks((0, 250, 500, 750, 1000))
    ax.set_yticks((0, 1000, 2000, 3000, 4000))
    ax.set_title(rf"{divergence}\% divergence" + "\n" +
                 rf"{mean_totals[divergence]:,.0f} mean anchors", fontsize=8)
    ax.tick_params(axis="both", labelsize=7)
    ax.grid(axis="y", color="#e5e5e5", linewidth=0.45)
    ax.set_xlabel(r"DP score $f_i$", fontsize=8)
axes[0].set_ylabel("mean anchors at exact score", fontsize=8)
axes[0].text(48, 4100, r"$\mathtt{min\_sc}=40$", rotation=90,
             va="top", ha="left", color="#A95700", fontsize=6.5)
fig.subplots_adjust(left=0.085, right=0.995, bottom=0.18, top=0.82, wspace=0.10)

print(get_filename_from_extension(".pgf"))
plt.savefig(get_filename_from_extension(".pgf"), bbox_inches="tight")
\end{CacheMeCode}
\caption{Exact initial-DP anchor score counts when the target contains
\(N=4000\) planted occurrences of the query source.  Here \(f_i\) is the DP
score of the best chain ending at anchor \(i\).  Each point is the arithmetic
mean of the six input-level counts at that integer score, including zero
counts in the mean.  Scores absent from all six inputs are blank.  The panels
show 0, 5, and 10 percent requested query substitutions.  No score-width binning or
smoothing is applied.  The dashed line is the default candidate-end threshold
\(\texttt{min\_sc}=40\).}
\label{fig:chaining-max-drop-score-counts}
\end{figure}

The \(N=4000\) panels contain approximately twenty times as many anchors as
the corresponding \(N=200\) panels, while their occupied score values and
shapes are nearly the same.  Across all 41 divergence settings, the ratio of
the two median anchor counts stayed between 19.9938 and 20.0000.  The largest
difference in the median percentage at or above score 40 was only 0.0157
percentage point.  As \(N\) increased in this construction, it therefore
scaled the amount of data but did not materially change the score-distribution
shape.

Figure~\ref{fig:chaining-max-drop-score-divergence} follows the \(N=4000\)
distribution over all 41 divergence values.  Score percentiles are first
calculated within each generated input using the lower-rank empirical
quantile, which is the first score whose cumulative count reaches the
requested fraction of anchors.  The solid curves in the upper panel are
across-input medians of those percentiles.  Its nested bands describe
within-input score intervals.  The lower panel shows the across-input median
fraction reaching \(\texttt{min\_sc}\), with the between-input interquartile
range.

\begin{figure}[H]
\centering
\begin{CacheMeCode}{python, custom include command={\input{\robExtAddCachePathAndName{\robExtFinalHash.pgf}}}, set placeholder eval={__LINEWIDTH__}{\lenToCmNoUnit[in]{\linewidth}}}
import matplotlib
matplotlib.use("pgf")
import matplotlib.pyplot as plt

matplotlib.rcParams.update({
    "font.family": "serif",
    "font.serif": [],
    "text.usetex": True,
})

# divergence, q10, q25, q50, q75, q90, fraction>=40 median, q1, q3
score_data = [
    (0.00, 109.5, 258.0, 501.5, 745.5, 894.5, 96.89655379, 96.63857678, 97.66870700),
    (0.25, 107.5, 246.0, 485.0, 725.5, 868.5, 97.35982328, 96.63312802, 97.70091896),
    (0.50, 100.0, 242.0, 473.0, 707.5, 846.5, 97.55349580, 96.67565969, 97.62553895),
    (0.75, 101.0, 230.5, 469.0, 695.5, 833.5, 97.19355847, 96.42530488, 97.49562322),
    (1.00, 101.0, 231.0, 454.0, 675.0, 821.0, 97.32294958, 96.89571365, 97.41092243),
    (1.25, 94.5, 235.0, 447.5, 661.5, 800.5, 97.25004953, 96.76691552, 97.33757050),
    (1.50, 94.0, 232.5, 435.5, 645.5, 784.5, 97.07967219, 96.61841419, 97.21161132),
    (1.75, 95.5, 239.0, 426.5, 628.5, 762.0, 97.12096376, 96.99197996, 97.61667017),
    (2.00, 87.5, 224.0, 411.5, 613.5, 741.0, 96.87481226, 96.37168933, 97.00258548),
    (2.25, 88.5, 224.0, 404.0, 599.0, 723.0, 96.85883222, 96.11322138, 97.36742424),
    (2.50, 87.0, 211.0, 395.0, 581.5, 699.5, 96.72547447, 95.95996493, 97.28515103),
    (2.75, 85.5, 202.0, 390.0, 567.5, 683.0, 96.68606642, 95.81185071, 97.14837824),
    (3.00, 83.0, 196.5, 378.0, 560.5, 672.0, 96.79292929, 95.68181818, 97.33187567),
    (3.25, 91.0, 192.0, 372.0, 556.5, 659.0, 96.67505392, 96.09728507, 97.29014133),
    (3.50, 91.0, 195.0, 363.5, 533.5, 640.0, 97.12912428, 96.30854912, 97.32903546),
    (3.75, 87.0, 191.5, 344.0, 515.5, 608.0, 96.96939151, 96.23469388, 97.24091038),
    (4.00, 88.5, 187.5, 342.0, 505.0, 601.0, 96.87364698, 96.10934415, 97.44898154),
    (4.25, 82.5, 178.5, 334.0, 497.0, 593.5, 96.82317745, 95.99329325, 97.35180612),
    (4.50, 80.5, 175.5, 329.0, 475.0, 574.5, 96.75590791, 95.72754831, 97.27905286),
    (4.75, 70.5, 169.5, 317.0, 455.0, 553.5, 95.87486157, 94.69993477, 97.12790698),
    (5.00, 69.0, 159.5, 303.0, 431.5, 531.5, 95.49006838, 94.42329654, 97.06935852),
    (5.25, 67.5, 154.0, 292.0, 412.5, 510.5, 94.32595573, 94.26108374, 96.50821596),
    (5.50, 61.5, 148.5, 280.5, 398.5, 498.5, 94.20168067, 93.98395722, 96.42857143),
    (5.75, 59.5, 141.5, 277.5, 402.5, 483.0, 93.67088608, 93.65581676, 96.03676914),
    (6.00, 59.5, 141.5, 270.0, 393.0, 471.5, 93.51618327, 93.44262295, 96.01648352),
    (6.25, 54.5, 129.5, 255.0, 385.5, 465.5, 93.27683616, 93.16809844, 95.83333333),
    (6.50, 54.5, 126.0, 249.0, 359.5, 451.0, 92.95774648, 92.88229376, 95.73943662),
    (6.75, 56.0, 120.5, 239.5, 361.0, 433.5, 93.43563512, 92.52802844, 95.89783282),
    (7.00, 56.0, 124.0, 236.0, 359.5, 431.5, 93.93382353, 92.84048507, 95.75163399),
    (7.25, 58.0, 124.0, 229.0, 337.0, 420.0, 93.60544218, 92.56410256, 95.58477237),
    (7.50, 58.0, 112.5, 221.5, 317.0, 382.5, 93.29085457, 92.02586207, 95.36956522),
    (7.75, 54.0, 104.0, 207.5, 295.5, 373.5, 92.27433977, 90.73849879, 93.08340016),
    (8.00, 51.5, 104.0, 206.5, 291.5, 362.5, 91.81818182, 90.07421150, 92.82467532),
    (8.25, 51.5, 109.0, 204.5, 290.0, 356.0, 91.60839161, 90.07421150, 92.58911820),
    (8.50, 44.0, 104.0, 196.5, 274.5, 346.0, 91.11505122, 89.53301127, 92.38453555),
    (8.75, 44.0, 100.5, 188.5, 269.0, 339.5, 91.02254428, 89.35185185, 91.74500588),
    (9.00, 44.0, 93.0, 187.0, 246.5, 321.0, 90.43040293, 87.59615385, 91.53796654),
    (9.25, 41.0, 94.0, 184.5, 248.0, 313.0, 90.13605442, 86.98386331, 91.19047619),
    (9.50, 35.5, 94.5, 178.5, 241.0, 295.5, 89.89795918, 89.55424275, 91.87500000),
    (9.75, 35.5, 91.5, 169.0, 225.0, 279.0, 89.41769317, 89.30471125, 91.28733997),
    (10.0, 35.5, 87.0, 154.0, 215.0, 260.5, 89.39280360, 88.62876254, 91.16379310),
]

(divergence, q10, q25, q50, q75, q90,
 fraction, fraction_q1, fraction_q3) = map(list, zip(*score_data))

fig, axes = plt.subplots(2, 1, sharex=True)
fig.set_size_inches(0.98 * __LINEWIDTH__, 0.53 * __LINEWIDTH__, forward=True)

axes[0].fill_between(divergence, q10, q90, color="#4C78A8", alpha=0.14,
                     label="10th--90th percentile")
axes[0].fill_between(divergence, q25, q75, color="#4C78A8", alpha=0.30,
                     label="25th--75th percentile")
axes[0].plot(divergence, q50, color="#285F8F", linewidth=1.3,
             label="median score")
axes[0].axhline(40, color="#F28E2B", linestyle="--", linewidth=0.9,
                label=r"$\mathtt{min\_sc}=40$")
axes[0].set_ylim(0, 1000)
axes[0].set_ylabel(r"DP score $f_i$", fontsize=8)
axes[0].set_title("(a) Within-input score spread", fontsize=8, loc="left")
axes[0].legend(frameon=False, fontsize=6.5, ncol=2, loc="upper right")

axes[1].fill_between(divergence, fraction_q1, fraction_q3,
                     color="#59A14F", alpha=0.22,
                     label="six-input interquartile range")
axes[1].plot(divergence, fraction, color="#3C7C36", linewidth=1.3,
             label="six-input median")
axes[1].set_ylim(85, 100)
axes[1].set_yticks((85, 90, 95, 100))
axes[1].set_ylabel(r"anchors with $f_i\geq40$ (\%)", fontsize=8)
axes[1].set_xlabel(r"query substitutions requested (\%)", fontsize=8)
axes[1].set_title(r"(b) Anchors reaching $\mathtt{min\_sc}$", fontsize=8, loc="left")
axes[1].legend(frameon=False, fontsize=6.5, loc="lower left")

for ax in axes:
    ax.set_xlim(0, 10)
    ax.set_xticks((0, 2, 4, 6, 8, 10))
    ax.tick_params(axis="both", labelsize=7)
    ax.grid(axis="y", color="#e5e5e5", linewidth=0.45)
fig.subplots_adjust(left=0.10, right=0.985, bottom=0.12, top=0.97, hspace=0.23)

print(get_filename_from_extension(".pgf"))
plt.savefig(get_filename_from_extension(".pgf"), bbox_inches="tight")
\end{CacheMeCode}
\caption{DP-score distribution across all 41 divergence settings when the
target contains \(N=4000\) planted occurrences.  The score \(f_i\) belongs to
the best chain ending at anchor \(i\).  The upper panel shows across-input
medians of within-input score percentiles.  The nested bands therefore describe
score spread, not uncertainty.  The lower panel shows the across-input median
percentage of anchors reaching the default \(\texttt{min\_sc}=40\) candidate-end
threshold, with the six-input interquartile range.  All ranges are descriptive.}
\label{fig:chaining-max-drop-score-divergence}
\end{figure}

At 0, 5, and 10 percent divergence, the across-input median DP
scores were 501.5, 303, and 154.  The corresponding median percentages with
\(f_i\geq40\) were 96.90, 95.49, and 89.39 percent, and the median DP-anchor
counts were 716,000, 314,000, and 128,000.  The median score and the
90th percentile decreased at every quarter-point divergence step.  The
median 10th percentile first fell below \texttt{min\_sc} at 9.5 percent
divergence.  The small local rebounds in the fraction above 40 mean that this
fraction should not be described as strictly monotone.

These marginal distributions explain the scale and range of \(f_i\), but
they do not by themselves predict \texttt{max\_drop}.  The drop rule depends
on the ordered score changes along predecessor paths, rather than on the
unordered histogram of all anchor scores.

\subsection{Default \texttt{max\_drop} on ordinary repeated copies}

The occurrence-by-divergence screen found no local opportunity at minimap2's
default \(D=500\).  In the 4.5-Mbp primary arm, none of 2,869,628 executed
searches crossed the strict threshold.  The fixed-start scan examined all
255,965,133 baseline scan steps, leaving no suffix.  The separate
\(N=5000\) sensitivity arm required a 5.625-Mbp reference and gave the same
result.  Across both feasible arms, none of 4,099,647 searches crossed
\(D=500\), and the fixed-start calculation retained all 365,680,152 baseline scan
steps.  No \(D=500\) cutoff fired.  The audited complete replay nevertheless
had the same aggregate search and scan-step counts and the same accepted-chain
output as the baseline.  Its execution trace and final marking nevertheless
differed in 57 of 4,920 observations because minimap2's best-cut boundary can
differ from the baseline terminal without a threshold crossing.  The matched
no-drop replay produced the same 57 differences, which isolates them from
the \(D=500\) rule.

A threshold sensitivity used
\(D\in\{0,25,50,100,200,500,750,1000\}\).  Thresholds 0 and 25 each fired on
57 of 4,099,647 searches, or 0.001390 percent, but each crossing occurred on
the final scan step and left no suffix.  No threshold of 50 or greater fired.  The
largest observed baseline-path drawdown was 39.  The SIMD replay agreed: at
each width \(L\in\{4,8,16\}\), \(D=0\) and \(D=25\) each produced 57 batches
containing a crossing, but no lane scan step or batch-critical step was removed.
No SIMD batch crossed \(D\geq50\).

The first screen kept the query length near 1,000 bases.  A second operational
screen therefore crossed target lengths 1.125, 4.5, and 18~Mbp with query
lengths near 250, 1,000, 4,000, and 16,000 bases.  A separate 72-Mbp/16-kbp
row extends a four-point constant-density diagonal.  On this diagonal, the
combined length of the five source-template occurrences is
\(5Q/T\approx0.111\) percent, where \(Q\) is query length and \(T\) is
target length.  This holds for the
1.125-Mbp/250-bp, 4.5-Mbp/1-kbp, 18-Mbp/4-kbp, and 72-Mbp/16-kbp cases.
Every configuration used five target occurrences of the unmutated source
template, the same six fixed sequence seeds, and all 41 requested substitution
levels from 0 to 10 percent in 0.25-percentage-point steps.  Thus each row in
Table~\ref{tab:chaining-max-drop-length-screen} contains 246 observations.

Five copies were used to preserve minimap2's default occurrence filtering.
In a predeclared \(N=50\) feasibility pilot, the default filter removed all
50-copy minimizers from the sparse 18-Mbp/247-base case and produced no
mapping.  That failed pilot was excluded before the final acquisition.  No
filter parameter was relaxed.  In the final screen every observation had a
forward PAF placement covering at least half of the query and half of a
planted locus.

The table uses the following labels.  \(V\) is the count of observations
passing that mapping gate.  \(E\) is the count of executed baseline searches.
\(H_{500}\) counts strict \(D=500\) crossings, while \(P_{500}\)
counts productive crossings with a nonempty suffix.  \(B\) is the baseline
backtracking scan-step count.  \(S_{500}\) counts steps after the first crossing
when the start and incoming marks are held fixed.  For each executed search,
\(\Delta(e)=\max_m d_m\) is measured while following its predecessor links.
The table reports \(\Delta_{\max}=\max_e\Delta(e)\), the largest observed
value in the row.  All are raw logical counters, not elapsed times.

\begin{table}[H]
\centering
\footnotesize
\setlength{\tabcolsep}{2.5pt}
\renewcommand{\arraystretch}{1.05}
\begin{tabular}{@{}rrrrrrrrr@{}}
\toprule
\(T\) (Mbp) & query length (bp) & \(V/246\) & \(E\) & \(H_{500}\) &
\(P_{500}\) & \(B\) & \(S_{500}\) & \(\Delta_{\max}\) \\
\midrule
1.125  & 247--252     & 246/246 & 1\,230 & 0 & 0 & 25\,705    & 0 & 0 \\
1.125  & 997--1002    & 246/246 & 1\,230 & 0 & 0 & 109\,715   & 0 & 0 \\
1.125  & 3997--4002   & 246/246 & 1\,230 & 0 & 0 & 440\,525   & 0 & 0 \\
1.125  & 15997--16002 & 246/246 & 1\,230 & 0 & 0 & 1\,757\,774 & 0 & 0 \\
\midrule
4.500  & 247--252     & 246/246 & 1\,230 & 0 & 0 & 25\,705    & 0 & 0 \\
4.500  & 997--1002    & 246/246 & 1\,230 & 0 & 0 & 109\,715   & 0 & 0 \\
4.500  & 3997--4002   & 246/246 & 1\,230 & 0 & 0 & 440\,525   & 0 & 0 \\
4.500  & 15997--16002 & 246/246 & 1\,230 & 0 & 0 & 1\,757\,741 & 0 & 0 \\
\midrule
18.000 & 247--252     & 246/246 & 1\,230 & 0 & 0 & 25\,670    & 0 & 0 \\
18.000 & 997--1002    & 246/246 & 1\,230 & 0 & 0 & 109\,715   & 0 & 0 \\
18.000 & 3997--4002   & 246/246 & 1\,230 & 0 & 0 & 440\,525   & 0 & 0 \\
18.000 & 15997--16002 & 246/246 & 1\,230 & 0 & 0 & 1\,757\,741 & 0 & 0 \\
\midrule
72.000\(^*\) & 15997--16002 & 246/246 & 1\,230 & 0 & 0 & 1\,757\,675 & 0 & 0 \\
\bottomrule
\end{tabular}
\caption{Default-\(D=500\) target- and query-length screen.  Each row combines
six seed blocks and 41 requested substitution percentages.  Count columns
are sums over those observations, while \(\Delta_{\max}\) is their maximum.
\(V\) is valid mapped
observations, \(E\) is baseline scans, \(H_{500}\) is threshold crossings,
\(P_{500}\) is productive crossings, \(B\) is baseline scan steps,
\(S_{500}\) is the count of remaining steps after the first crossing, and
\(\Delta_{\max}\) is the largest score drop.  The starred 72-Mbp row extends the constant-density diagonal and
is not part of the 3-by-4 screen.}
\label{tab:chaining-max-drop-length-screen}
\end{table}

All 3,198 observations mapped, giving 15,990 executed searches and 8,758,731
baseline scan steps.  None of the \(13\times41=533\)
configuration-by-divergence cells contained a strict \(D=500\) crossing.
This zero occurs in every cell, not only after the counts are added.  The audited
complete max-drop replay also made 15,990 searches and 8,758,731 scan steps.  Its
execution trace, final marks, and accepted-chain output matched the baseline
in all 3,198 observations.  The approximately 250-base
tier is a structural boundary case.  A chain score is bounded by the query
span credited along that segment, which is at most 252 bases here, so its
score range cannot produce a drawdown above 500.  Its 41 requested substitution
levels also collapse to 26 distinct rounded mutation counts.  The longer
queries are the informative rows.  Their maximum drawdown was nevertheless
also zero.

The zero values in the last column are measured values, not rounded values or
missing data.  The matched no-drop backtrack evaluated 8,758,731 values of
\(s_m\) while traversing the same paths as the separate drop diagnostic.  Every
value was larger than the preceding running maximum, so \(M_m=s_m\) and
\(d_m=0\) at every evaluated step.  This is possible even on a long chain.
These inputs contain collinear, substitution-only copies, and on every
predecessor edge traversed by those executed searches, the DP score increased
by a positive amount toward the chain end.  Moving backward therefore lowers
the prefix score \(f_{b_m}\), which makes \(s_m=f_e-f_{b_m}\) increase.  The
same code measured nonzero drops in the ordinary occurrence screen and in the
long-chain example below, so the zero is specific to this length screen.

Target length, planted-copy fraction, retained-anchor load, and adaptive
occurrence filtering change together in this construction.  The table is
therefore an operational screen, not a causal estimate of a target-length or
query-length effect.  Within this synthetic construction, no tested
target/query-length configuration produced a crossing.

\subsection{A long-chain example that crosses
\texorpdfstring{\(D=500\)}{D=500}}

A specially prepared long-chain example tests whether the instrumented
predicate can fire at all.  The sequence letters are pseudorandom, but their
layout is deliberately chosen.  The 12,480-base query is a 6,000-base left
flank, followed by 16 adjacent 30-base exact pieces, followed by a 6,000-base
right flank.  The 19,980-base target contains the same parts in the same
order, but inserts a 500-base target-only gap between every neighboring pair
of exact pieces.  There are therefore 15 inserted gaps.  Figure~\ref{fig:chaining-max-drop-long-chain-accounting}(a)
shows this construction.

The horizontal variable in panels (b) and (c) is the requested substitution
percentage \(x_k\), defined by
\[
  x_k=0.25k,
  \qquad
  m_k=\left\lfloor
    12{,}480\frac{x_k}{100}+\frac{1}{2}
  \right\rfloor,
  \qquad k=0,\ldots,40.
\]
If \(u_{k,r}=1\) when query base \(r\) is changed at level \(k\), and zero
otherwise, then \(m_k=\sum_{r=1}^{12{,}480}u_{k,r}\).  In simple terms,
\(m_k\) is the number of query bases that are changed.  These positions
are chosen only from the 12,000 flank bases, and the sets are nested as
\(k\) increases.  All 480 bases in the 16 exact pieces remain unchanged.  For
example, \(x_k=5\) changes \(m_k=624\) query bases, all in the two flanks.
Thus \(x_k\) uses the complete 12,480-base query as its denominator, but it
does not mean that five percent of every region is changed.

For panel (b), each baseline scan keeps its original starting end and incoming
marks.  The calculation finds the first step with \(d_m>500\) and counts the
remaining baseline steps.  It does not update the marks.  Panel (c) instead
applies every cutoff and updates the marks before the next scan.  At each
\(x_k\), each plotted percentage is a ratio of sums: its numerator and
denominator are first added across the six fixed sequence seeds.  The
16-piece layout was selected after exploratory 8- and 12-piece examples did
not cross \(D=500\).  It demonstrates that the rule can fire, but it does not
estimate how often this happens in ordinary inputs.

\begin{figure}[H]
\centering
\begin{CacheMeCode}{python, custom include command={\input{\robExtAddCachePathAndName{\robExtFinalHash.pgf}}}, set placeholder eval={__LINEWIDTH__}{\lenToCmNoUnit[in]{\linewidth}}}
import matplotlib
matplotlib.use("pgf")
import matplotlib.pyplot as plt

matplotlib.rcParams.update({
    "font.family": "serif",
    "font.serif": [],
    "text.usetex": True,
})

from matplotlib.patches import Rectangle

# Requested substitution percentage, the fraction of one fixed-start scan
# after its first crossing, and the complete-replay scan-step change.
long_chain_data = [
    (0.00,49.50771492,1.95444526),(0.25,49.54381661,2.03940811),
    (0.50,49.40009598,2.12765957),(0.75,49.45421215,2.21648196),
    (1.00,49.62654160,2.31023102),(1.25,49.29270947,2.41204207),
    (1.50,49.08352230,2.51322751),(1.75,49.28895912,2.62690105),
    (2.00,49.26516249,2.75305320),(2.25,49.28633218,2.87629758),
    (2.50,49.67283394,3.00090253),(2.75,49.84097067,3.13346684),
    (3.00,49.76687117,3.26380368),(3.25,49.71187092,3.40632603),
    (3.50,49.85938128,3.56234097),(3.75,49.63035291,3.71041986),
    (4.00,49.47598253,3.87190684),(4.25,48.96341463,4.05487805),
    (4.50,49.03493380,4.24310097),(4.75,49.25422605,4.40835267),
    (5.00,49.42251336,4.58541631),(5.25,49.63950973,4.79452055),
    (5.50,50.15054573,5.00564546),(5.75,50.29492725,5.23004326),
    (6.00,50.23464599,5.42746378),(6.25,50.08481764,5.64037320),
    (6.50,49.80106101,5.87975243),(6.75,50.06938020,6.15171138),
    (7.00,50.20546290,6.42978003),(7.25,50.07571933,6.71378092),
    (7.50,50.02626050,6.98529412),(7.75,50.24711697,7.30367930),
    (8.00,50.15754798,7.61959324),(8.25,50.20908005,7.94504182),
    (8.50,49.56222639,8.31769856),(8.75,49.43200260,8.63356053),
    (9.00,49.39597315,8.92617450),(9.25,48.65054329,9.32351910),
    (9.50,48.52027768,9.71867008),(9.75,48.41057066,10.18766756),
    (10.00,48.53057983,10.56393963),
]
substitution_percent, remaining_percent, full_extra_percent = map(
    list, zip(*long_chain_data)
)

fig = plt.figure()
fig.set_size_inches(0.91 * __LINEWIDTH__, 0.68 * __LINEWIDTH__, forward=True)
grid = fig.add_gridspec(2, 2, height_ratios=(1.0, 1.18), hspace=0.56,
                        wspace=0.27)
top = fig.add_subplot(grid[0, :])
left = fig.add_subplot(grid[1, 0])
right = fig.add_subplot(grid[1, 1])

flank_color = "#B8C7D9"
piece_color = "#2F6B9A"
gap_color = "#E17C05"
edge_color = "#3D4A57"
x_start = 0.055
flank_width = 0.160
piece_width = 0.008
gap_width = 0.018
block_height = 0.22
query_y = 1.25
target_y = 0.48

top.set_xlim(0, 1)
top.set_ylim(0, 2.02)
top.axis("off")
top.set_title("(a) Construction of the long-chain example (not to scale)",
              fontsize=8, loc="left")

def add_block(x, y, width, color, edge=edge_color, linewidth=0.45):
    top.add_patch(Rectangle((x, y), width, block_height,
                            facecolor=color, edgecolor=edge,
                            linewidth=linewidth))

def add_bracket(x_left, x_right, y, text_value, text_y):
    top.plot([x_left, x_left, x_right, x_right],
             [y - 0.035, y, y, y - 0.035],
             color=edge_color, linewidth=0.55)
    top.text((x_left + x_right) / 2, text_y, text_value,
             ha="center", va="center", fontsize=6.3)

query_middle_start = x_start + flank_width
query_middle_end = query_middle_start + 16 * piece_width
query_end = query_middle_end + flank_width
target_middle_start = query_middle_start
target_middle_end = target_middle_start + 16 * piece_width + 15 * gap_width
target_end = target_middle_end + flank_width

add_block(x_start, query_y, flank_width, flank_color)
for piece_index in range(16):
    add_block(query_middle_start + piece_index * piece_width, query_y,
              piece_width, piece_color, edge="white", linewidth=0.30)
add_block(query_middle_end, query_y, flank_width, flank_color)

add_block(x_start, target_y, flank_width, flank_color)
target_x = target_middle_start
for piece_index in range(16):
    add_block(target_x, target_y, piece_width, piece_color,
              edge="white", linewidth=0.30)
    target_x += piece_width
    if piece_index < 15:
        add_block(target_x, target_y, gap_width, gap_color,
                  edge="white", linewidth=0.30)
        target_x += gap_width
add_block(target_middle_end, target_y, flank_width, flank_color)

top.text(0.045, query_y + block_height / 2, "query", ha="right",
         va="center", fontsize=7)
top.text(0.045, target_y + block_height / 2, "target", ha="right",
         va="center", fontsize=7)
for x_position, y_position in (
    (x_start + flank_width / 2, query_y),
    (query_middle_end + flank_width / 2, query_y),
    (x_start + flank_width / 2, target_y),
    (target_middle_end + flank_width / 2, target_y),
):
    top.text(x_position, y_position + block_height / 2,
             "6,000-bp\nflank", ha="center", va="center", fontsize=5.7)

add_bracket(query_middle_start, query_middle_end,
            query_y + block_height + 0.07,
            "16 adjacent exact pieces, 30 bp each",
            query_y + block_height + 0.14)
add_bracket(target_middle_start, target_middle_end,
            target_y - 0.07,
            "15 target-only gaps, 500 bp each",
            target_y - 0.16)
top.text((x_start + query_end) / 2, 1.05,
         r"At $x_k>0$, only query-flank bases are changed.",
         ha="center", va="center", fontsize=6.2, color="#8A3E00")
top.text(0.80, 1.36,
         r"$|Q|=6000+16(30)+6000$" "\n" r"$=12{,}480$ bp",
         ha="left", va="center", fontsize=6.4)
top.text(0.80, 0.59,
         r"$|T|=|Q|+15(500)$" "\n" r"$=19{,}980$ bp",
         ha="left", va="center", fontsize=6.4)

left.plot(substitution_percent, remaining_percent,
          color="#4C78A8", linewidth=1.25)
left.axhline(49.54799892, color="#7F7F7F", linestyle="--", linewidth=0.8,
             zorder=1)
left.set_xlim(0, 10)
left.set_ylim(47.5, 51.0)
left.set_xticks((0, 2.5, 5, 7.5, 10))
left.set_yticks((48, 49, 50, 51))
left.set_title("(b) Stop one scan at its first crossing", fontsize=8, loc="left")
left.set_ylabel("steps after first crossing\n" r"/ baseline steps (\%)",
                fontsize=8)
left.set_xlabel(r"requested substitutions $x_k$ (\%)", fontsize=8)
left.text(0.35, 50.65,
          "all 246 observations\n" r"combined: 49.548\%",
          fontsize=6.5, color="#555555", va="top")

right.plot(substitution_percent, full_extra_percent,
           color="#D55E00", linewidth=1.25)
right.axhline(0, color="#555555", linewidth=0.7)
right.set_xlim(0, 10)
right.set_ylim(0, 12)
right.set_xticks((0, 2.5, 5, 7.5, 10))
right.set_yticks((0, 3, 6, 9, 12))
right.set_title("(c) Apply every cutoff and update marks", fontsize=8, loc="left")
right.set_ylabel(r"extra scan steps vs. baseline (\%)", fontsize=8)
right.set_xlabel(r"requested substitutions $x_k$ (\%)", fontsize=8)
right.text(0.35, 10.8,
           "all 246 observations\n" r"combined: +4.046\%",
           fontsize=6.5, color="#A34700", va="top")

for ax in (left, right):
    ax.tick_params(axis="both", labelsize=7)
    ax.grid(axis="y", color="#e5e5e5", linewidth=0.45)
fig.subplots_adjust(left=0.14, right=0.985, bottom=0.10, top=0.98)

print(get_filename_from_extension(".pgf"))
plt.savefig(get_filename_from_extension(".pgf"), bbox_inches="tight")
\end{CacheMeCode}
\caption{Construction and scan-step accounting for the long-chain example at
default \(D=500\).  Panel (a) shows the specially prepared layout.  Its
sequence letters are pseudorandom.  The 12,480-base query has two 6,000-base
flanks around 16 adjacent exact 30-base pieces.  The 19,980-base target adds
one 500-base target-only gap between each neighboring pair.  The drawing is
not to scale.  At \(x_k>0\), substitutions are made only in the query flanks.
Panel (b) counts the part of each fixed-start baseline scan after its first
strict crossing.  The dashed value is \(\sum S/\sum B\) over all 246
observations.  Panel (c) applies every cutoff and updates the marks, which
increases the complete scan work.  These are logical scan steps, not
instructions, elapsed time, or a measured slowdown.}
\label{fig:chaining-max-drop-long-chain-accounting}
\end{figure}

All 246 long-chain observations retained one long-span, plus-strand, MAPQ-60
placement on the sole intended target, with at least 98.910 percent query-axis
coverage.  This validates the intended chaining placement, not a high-identity
alignment: the residue-match count divided by query length fell to 26.402
percent, and only 147 of 246 observations reached 50 percent.

Holding each baseline start and its incoming marks fixed, all 246 paths
crossed productively.  These scans reached the first crossing after 136,008 of 269,579
baseline scan steps,
leaving a local suffix of 133,571 steps, or 49.548 percent.  Its
\(\sum S/C=0.9821\) is a large local opportunity under
Equation~\eqref{eq:chaining-max-drop-break-even}.  It does not compose into a
complete-work reduction.  The complete replay left suffix anchors unmarked,
then processed the newly exposed candidate ends.  It made 7,503
searches and 280,485 scan steps, compared with 246 searches and
269,579 steps for the baseline policy: 10,906 more steps, or 4.046 percent,
before charging any added check cost.  It accepted 492 chains and rejected
7,011 paths, whereas the baseline policy accepted 246 and rejected none.  The
accepted-chain output of this initial-DP replay changed in all 246
observations.  This is an internal chaining result, not a comparison of final
PAF records.  The no-drop best-cut replay
kept the same bookkeeping but disabled the score-drop cutoff and matched
the baseline policy exactly.  The difference therefore came from the
\(D=500\) cutoff decisions rather than from that independent bookkeeping rule.

The fraction of steps after the first crossing stayed near one half across the
substitution sweep, but the complete-work penalty grew from 1.954 percent at
\(x_k=0\) to 10.564 percent at \(x_k=10\).  The absolute excess was 266 scan steps over the six seeds
at every level.  The percentage rose because baseline path work decreased.

For both ordinary screens, \(\sum S/C=0\).  Even a free check can only tie the
baseline scan in the fixed-start variable-work model, while any positive check cost
makes that model worse.  The long-chain example proves that the predicate
can fire, but it increases complete logical scan work and changes chain
semantics.  These results therefore do not justify an always-on SIMD
\texttt{max\_drop} check for the studied BioseqDB path.  They do not prove the
rule useless for indels, noisy reads, other presets, or other reference
structures, and they do not quantify wall-clock performance.  A timing study
of baseline and checked BioseqDB kernels would still be required.

This max-drop conclusion does not predict \(\operatorname{Run}_L\), which
tests structural adjacency rather than score drawdown.  The next subsection
evaluates that separate existing fast-path condition.

\subsection{Frequency of the \texorpdfstring{\(\operatorname{Run}_L\)}{Run-L}
condition}

Figure~\ref{fig:chaining-max-drop-runl} reports the condition-true ratio
\(100B_L/E_L\), meaning the percentage of loop checks for which
\(\operatorname{Run}_L\) is true, and the ideal loop-iteration reduction
\(100(q-E_L)/q\).  Within each generated input, the measured counters were
identical for every \(N=2,\ldots,4000\), so one curve represents all 18 of
those \(N\) values.  Hollow markers
show only the divergence settings at which the six-input \(N=1\) median
differed.

\begin{figure}[H]
\centering
\begin{CacheMeCode}{python, custom include command={\input{\robExtAddCachePathAndName{\robExtFinalHash.pgf}}}, set placeholder eval={__LINEWIDTH__}{\lenToCmNoUnit[in]{\linewidth}}}
import matplotlib
matplotlib.use("pgf")
import matplotlib.pyplot as plt
from matplotlib.lines import Line2D

matplotlib.rcParams.update({
    "font.family": "serif",
    "font.serif": [],
    "text.usetex": True,
})

# divergence, Fire_4, Fire_8, Fire_16, Reduction_4, Reduction_8, Reduction_16
# Each row is the six-input median common to every tested N>=2.
runl_data = [
    (0.00, 95.74275362, 88.14102564, 62.85714286, 74.17542588, 86.03620474, 90.18817204),
    (0.25, 93.40579710, 84.49074074, 48.80952381, 73.69924812, 85.52389260, 87.97729618),
    (0.50, 93.25075834, 83.33333333, 66.96428571, 73.66695881, 85.36585366, 90.91243021),
    (0.75, 90.80338266, 82.60869565, 49.20634921, 73.14786585, 85.25641026, 87.91208791),
    (1.00, 96.05263158, 88.18181818, 50.00000000, 74.23462352, 86.05391074, 88.23529412),
    (1.25, 90.00000000, 71.41666667, 45.51630435, 72.97297297, 83.33024591, 86.66002871),
    (1.50, 93.24324324, 79.45652174, 42.10526316, 73.66247847, 84.70716385, 86.33093525),
    (1.75, 91.73745174, 73.32015810, 57.14285714, 73.33735733, 83.69245323, 89.55223881),
    (2.00, 91.42156863, 70.49689441, 43.33333333, 73.28083990, 83.14814815, 86.11111111),
    (2.25, 90.76704545, 75.71428571, 43.92156863, 73.13998770, 84.09090909, 86.78278689),
    (2.50, 91.93548387, 80.06535948, 63.63636364, 73.38614591, 84.85007496, 90.51724138),
    (2.75, 90.74952562, 74.34640523, 40.00000000, 73.13591236, 83.87274890, 85.41666667),
    (3.00, 90.66937120, 72.44582043, 35.08771930, 73.10850440, 83.49326599, 84.00537634),
    (3.25, 88.37597330, 66.66666667, 46.42857143, 72.60931242, 82.35294118, 87.38687783),
    (3.50, 93.23821340, 70.58823529, 48.70129870, 73.65045200, 83.16831683, 87.82368621),
    (3.75, 89.70223325, 75.00000000, 47.85714286, 72.89540816, 84.00000000, 87.13483146),
    (4.00, 89.33333333, 74.90196078, 44.44444444, 72.81323877, 83.97745572, 86.30952381),
    (4.25, 90.64285714, 65.68627451, 37.78280543, 73.10989465, 82.13391740, 84.45018995),
    (4.50, 89.23809524, 66.96832579, 28.78289474, 72.80219780, 82.40481523, 81.10240823),
    (4.75, 91.48550725, 73.21428571, 32.05128205, 73.29457364, 83.66666667, 82.76255708),
    (5.00, 88.47826087, 69.71153846, 37.41258741, 72.63106480, 82.86634912, 84.86715749),
    (5.25, 88.48684211, 57.98319328, 40.83333333, 72.63546798, 80.22988506, 85.96059113),
    (5.50, 89.18128655, 61.90476190, 40.00000000, 72.79220779, 81.17647059, 85.71428571),
    (5.75, 84.84848485, 51.92307692, 19.34984520, 71.79023508, 78.40501792, 74.26500612),
    (6.00, 88.09523810, 58.33333333, 20.48611111, 72.53846154, 80.32786885, 75.34678436),
    (6.25, 84.03361345, 66.81818182, 26.96969697, 71.59322034, 82.35875706, 80.17857143),
    (6.50, 80.99415205, 56.66666667, 36.36363636, 70.82089552, 79.82123510, 84.50704225),
    (6.75, 92.58241758, 68.33333333, 40.97222222, 73.52714671, 82.70189432, 85.93109106),
    (7.00, 82.10526316, 58.33333333, 28.20512821, 71.11501317, 80.06535948, 80.45977011),
    (7.25, 85.04273504, 74.67532468, 33.92857143, 71.73469388, 83.69047619, 82.74291498),
    (7.50, 83.97435897, 67.53246753, 21.53846154, 71.58385093, 82.50000000, 76.29310345),
    (7.75, 81.17647059, 68.33333333, 19.04761905, 70.88733799, 82.70189432, 73.84987893),
    (8.00, 81.25000000, 56.41025641, 30.00000000, 70.90909091, 79.35828877, 81.81818182),
    (8.25, 82.93269231, 63.33333333, 36.42857143, 71.32411067, 81.56108597, 84.17832168),
    (8.50, 85.64102564, 47.72727273, 12.91666667, 71.98067633, 76.93236715, 65.94202899),
    (8.75, 85.64102564, 47.72727273, 13.80952381, 71.98067633, 76.93236715, 67.42424242),
    (9.00, 77.35042735, 52.22222222, 18.33333333, 69.88372093, 78.22245322, 73.21428571),
    (9.25, 80.12820513, 59.82142857, 22.61904762, 70.59800664, 80.69767442, 76.25482625),
    (9.50, 79.79797980, 41.42857143, 22.50000000, 70.52631579, 74.34210526, 76.97368421),
    (9.75, 80.19480519, 42.85714286, 18.38235294, 70.63269877, 75.00000000, 71.38857783),
    (10.0, 79.80769231, 47.72727273, 10.79545455, 70.48494983, 76.93236715, 61.45484950),
]

# Actual N=1 medians are listed only where they differ from the common N>=2
# value.  Empty positions are deliberately not drawn.
n1_fire = {
    4: [(0.50,95.40169133),(0.75,100.0),(1.25,96.25),(1.75,97.29729730),(2.00,92.80303030),(2.25,92.32954545),(2.50,93.64919355),(2.75,96.49014778),(3.00,93.21839080),(3.25,91.60879630),(3.75,95.74275362),(4.00,98.0),(4.25,93.72727273),(4.50,91.23809524),(5.25,90.19138756),(5.50,89.47368421),(5.75,86.36363636),(6.50,84.60526316),(7.00,85.43859649),(7.25,92.82051282),(7.50,85.16483516),(8.25,87.76223776),(9.00,80.55555556),(9.50,85.90909091),(9.75,85.35353535),(10.0,86.05769231)],
    8: [(0.75,89.28571429),(1.25,76.90909091),(1.75,76.95652174),(2.75,79.41176471),(3.00,77.70897833),(3.25,70.83333333),(3.75,80.35714286),(4.25,70.0),(5.25,71.07843137),(5.75,55.04807692),(6.50,61.81818182),(7.00,69.69696970),(9.50,43.65079365)],
    16: [(0.75,63.49206349),(1.25,62.74038462),(3.25,50.0),(3.75,60.0),(4.25,48.07692308),(7.00,41.66666667)],
}
n1_reduction = {
    4: [(0.50,74.10701630),(0.75,75.0),(1.25,74.25496689),(1.75,74.47183099),(2.00,73.57142857),(2.25,73.47061176),(2.50,73.74965268),(2.75,74.32410490),(3.00,73.66028708),(3.25,73.31932773),(3.75,74.17542588),(4.00,74.61340206),(4.25,73.76095119),(4.50,73.24058920),(5.25,73.01393728),(5.50,72.85714286),(5.75,72.15189873),(6.50,71.73638848),(7.00,71.93200663),(7.25,73.57679914),(7.50,71.86956522),(8.25,72.45493107),(9.00,70.71428571),(9.50,72.01280228),(9.75,71.88995215),(10.0,72.07646177)],
    8: [(0.75,86.16118770),(1.25,84.28942187),(1.75,84.32565087),(2.75,84.73832528),(3.00,84.46833931),(3.25,83.17647059),(3.75,84.85714286),(4.25,83.02429668),(5.25,82.91281189),(5.75,79.38954675),(6.50,81.21794872),(7.00,82.96751536),(9.50,75.33783784)],
    16: [(0.75,90.37456446),(1.25,90.31009486),(3.25,88.23529412),(3.75,90.0),(4.25,87.80696745),(7.00,85.78431373)],
}

divergence = [row[0] for row in runl_data]
colors = {4: "#2F6B9A", 8: "#E17C05", 16: "#3C7C36"}
widths = {4: 1.55, 8: 1.05, 16: 1.05}

fig, axes = plt.subplots(2, 1, sharex=True)
fig.set_size_inches(0.98 * __LINEWIDTH__, 0.53 * __LINEWIDTH__, forward=True)
for column, lanes in enumerate((4, 8, 16), start=1):
    axes[0].plot(divergence, [row[column] for row in runl_data],
                 color=colors[lanes], linewidth=widths[lanes])
    axes[1].plot(divergence, [row[column + 3] for row in runl_data],
                 color=colors[lanes], linewidth=widths[lanes])
    for ax, exceptional in ((axes[0], n1_fire[lanes]),
                            (axes[1], n1_reduction[lanes])):
        if exceptional:
            x, y = zip(*exceptional)
            ax.scatter(x, y, s=13, facecolors="white", edgecolors=colors[lanes],
                       linewidths=0.75, zorder=4)

axes[0].set_title(r"(a) Checks where $\mathrm{Run}_L$ is true",
                  fontsize=8, loc="left")
axes[1].set_title(r"(b) Scalar loop iterations removed",
                  fontsize=8, loc="left")
axes[0].set_ylabel(r"condition true (\%)", fontsize=8)
axes[1].set_ylabel(r"ideal reduction (\%)", fontsize=8)
axes[1].set_xlabel(r"query substitutions requested (\%)", fontsize=8)
for ax in axes:
    ax.set_xlim(-0.1, 10.1)
    ax.set_ylim(0, 105)
    ax.set_xticks((0, 2, 4, 6, 8, 10))
    ax.set_yticks((0, 25, 50, 75, 100))
    ax.tick_params(axis="both", labelsize=7)
    ax.grid(axis="y", color="#e5e5e5", linewidth=0.45)

legend = [
    Line2D([0], [0], color=colors[4], linewidth=widths[4], label=r"$L=4$"),
    Line2D([0], [0], color=colors[8], linewidth=widths[8], label=r"$L=8$"),
    Line2D([0], [0], color=colors[16], linewidth=widths[16], label=r"$L=16$"),
    Line2D([0], [0], marker="o", markerfacecolor="white",
           markeredgecolor="#333333", linestyle="none", markersize=4,
           label=r"$N=1$ differs"),
]
fig.legend(handles=legend, frameon=False, fontsize=7, ncol=4,
           loc="lower center", bbox_to_anchor=(0.5, 0.005))
fig.subplots_adjust(left=0.09, right=0.985, bottom=0.19, top=0.98, hspace=0.25)

print(get_filename_from_extension(".pgf"))
plt.savefig(get_filename_from_extension(".pgf"), bbox_inches="tight")
\end{CacheMeCode}
\caption{BioseqDB \(\operatorname{Run}_L\) behavior along the highest-ranked
predecessor path in the ordinary synthetic repeated-copy inputs, not the
long-chain example.  Each of six fixed input blocks uses a pseudorandom
4.5-Mbp target background and one pseudorandom source of length 997--1002
bases.  For the 19 listed values from \(N=1\) to \(N=4000\), nonoverlapping
forward copies of the unchanged source were deliberately planted in the
target.  Queries were made from that source with nested substitutions from 0
to 10 percent.  Thus, at nonzero substitution percentages, \(N\) counts
homologous planted loci rather than exact copies of the changed query.  The
separate \(N=5000\) sensitivity arm is not shown.  Here \(L\) is the batch
width, \(B_L\) counts checks where \(\operatorname{Run}_L\) is true, \(E_L\)
counts all loop checks, and \(q\) is the scalar path length.  Solid curves are
six-input medians common to every tested \(N\geq2\).  Hollow markers show
where the \(N=1\) median differs.  The values are logical opportunities, not
measured speedups.}
\label{fig:chaining-max-drop-runl}
\end{figure}

After summing the numerators and denominators over all primary observations,
the firing ratios were 89.42, 70.05, and 37.15 percent for
\(L=4,8,16\).  The corresponding ideal loop-iteration reductions were 72.84,
83.06, and 84.79 percent.  At least one fast batch occurred in every one of
the 4,674 primary observations for every width.  This 100-percent incidence
does not mean that every evaluation was fast.  The upper panel gives that
denominator-specific ratio.

The median scalar path length fell from 179 anchors at zero divergence to
78.5 at 5 percent and 32 at 10 percent.  At the last point, the 16-wide
condition produced only one or two successful batches per input.  Its
percentage was consequently discrete and variable, with an observed firing
range of 6.25 to 66.67 percent.

Within each seed and divergence level, every \((q,B_L,E_L)\) counter was
identical for all 18 values \(N=2,\ldots,4000\).  Only \(N=1\) sometimes had
longer runs of consecutive predecessor indices.  Its largest median firing-rate differences were 9.20,
13.10, and 17.22 percentage points for \(L=4,8,16\), respectively.  Thus, as
\(N\) increased in this construction, the number of candidate ends increased
but no graded \(N\) trend appeared on the single highest-ranked path.  This
statement must not be extended to all later backtracking paths.

The four-wide case is the directly relevant Apple ARM64 width: four
32-bit predecessor indices fit in one 128-bit NEON vector.  It remained true
for a median 79.81 percent of loop evaluations even at 10 percent divergence,
with a 70.48-percent ideal reduction in loop evaluations.  The 8- and
16-wide results are what-if measurements for wider batches that require
multiple NEON registers or other scheduling.  They are not single-instruction
widths.  The evidence
therefore supports keeping the four-wide \(\operatorname{Run}_4\) fast path as
a frequent logical opportunity on this workload.  A BioseqDB timing benchmark
is still required to decide its actual gain.  The next section returns from
this empirical test to alternative chaining algorithms.



\clearpage
